% AxiomForge manuscript — template version 1.0.0.
% Build: latexmk -pdf main.tex   (see README.md)
\documentclass[11pt]{amsart}
\makeatletter
% AxiomForge manuscript preamble — template version 1.0.0.
% This file is copied into every contribution and is owned by it afterwards.
% It fixes the typographic identity of the series; it defines no mathematical notation
% (each paper keeps its own notation in macros.tex).

% ---------------------------------------------------------------- page and fonts
\usepackage[a4paper,hmargin=31mm,top=34mm,bottom=36mm,headsep=10mm,footskip=12mm]{geometry}
\usepackage[T1]{fontenc}
\usepackage[utf8]{inputenc}
\usepackage{lmodern}
\usepackage{amsmath,amssymb,amsthm,mathtools}
\usepackage[final,protrusion=true,expansion=true]{microtype}
\usepackage{booktabs}
\usepackage{array}
\usepackage[shortlabels]{enumitem}
\setlist{itemsep=0.2em,topsep=0.4em}
\usepackage{xcolor}
\linespread{1.04}
\setlength{\emergencystretch}{1.5em}
\allowdisplaybreaks[1]
\numberwithin{equation}{section}
\setcounter{tocdepth}{2}
% Table of contents: sections (and their page numbers) in bold with a little space above
% them, subsections indented in normal weight; after \appendix only the (bold) section
% titles are listed, so the appendices never crowd the contents.
\def\l@section{\@tocline{1}{2.5pt plus 1pt minus 2pt}{0pt}{}{\bfseries\let\@tocpagenum\af@tocpagenumbf}}
\def\af@tocpagenumbf#1{\hss{\bfseries #1}}
\def\l@subsection{\@tocline{2}{0pt}{1.5em}{2.3em}{}}
\g@addto@macro\appendix{\addtocontents{toc}{\protect\setcounter{tocdepth}{1}}}
% Every appendix starts on a new page; the unnumbered back matter (\afbacksection) keeps the
% original \section and flows normally after the last appendix.
% (Local implementation through a hook, so that \section itself is not redefined.)
\newif\ifaf@appendixpage
\AddToHook{cmd/section/before}{\ifaf@appendixpage\clearpage\fi}
\g@addto@macro\appendix{\global\af@appendixpagetrue}

% ---------------------------------------------------------------- links
\definecolor{afblue}{rgb}{0.08,0.18,0.42}
\usepackage[pdfusetitle,bookmarks=true,bookmarksnumbered=true,bookmarksopen=true,
  colorlinks=true,linkcolor=afblue,citecolor=afblue,urlcolor=afblue,
  pdfpagelabels=true,pdfdisplaydoctitle=true]{hyperref}
\usepackage[capitalise,noabbrev]{cleveref}
% Local addition: unique PDF anchors for theorem-like environments numbered within sections.
\hypersetup{hypertexnames=false}

% ---------------------------------------------------------------- theorem environments
% thmtools keeps cleveref names correct for environments sharing one counter.
\usepackage{thmtools}
\declaretheorem[style=plain,numberwithin=section,name=Theorem]{theorem}
\declaretheorem[style=plain,sibling=theorem,name=Proposition]{proposition}
\declaretheorem[style=plain,sibling=theorem,name=Lemma]{lemma}
\declaretheorem[style=plain,sibling=theorem,name=Corollary]{corollary}
\declaretheorem[style=plain,numbered=no,name=Theorem]{theorem*}
\declaretheorem[style=definition,sibling=theorem,name=Definition]{definition}
\declaretheorem[style=definition,sibling=theorem,name=Example]{example}
\declaretheorem[style=remark,sibling=theorem,name=Remark]{remark}
\crefname{equation}{}{}
\Crefname{equation}{Equation}{Equations}

% ---------------------------------------------------------------- series metadata
% \afidentifier{AF-0001}, \afversion{1.0.0}, \afdate{25 November 2026},
% \afrepository{https://github.com/axiomforge-math/mathematics/tree/main/results/...},
% \afdisclosure{short AI-assisted research disclosure},
% \afstatus{verification status} (local addition; optional)
\newcommand{\afidentifier}[1]{\gdef\af@identifier{#1}}
\newcommand{\afversion}[1]{\gdef\af@version{#1}}
\newcommand{\afdate}[1]{\gdef\af@date{#1}}
\newcommand{\afrepository}[1]{\gdef\af@repository{#1}}
\newcommand{\afdisclosure}[1]{\gdef\af@disclosure{#1}}
\newcommand{\afcorrespondence}[1]{\gdef\af@correspondence{#1}}
\gdef\af@correspondence{}
% Local addition: verification status of the contribution, printed as a 'Status' row.
\newcommand{\afstatus}[1]{\gdef\af@status{#1}}
\gdef\af@status{}

% The publication block printed under the abstract on the title page.
\newcommand{\afpublicationblock}{%
  \begingroup\footnotesize\noindent
  \begin{tabular}{@{}l@{\quad}p{0.72\textwidth}@{}}
    \textit{Identifier} & \af@identifier\quad(version \af@version, \af@date)\\[1pt]
    \ifx\af@status\@empty\else\textit{Status} & \af@status\\[1pt]\fi
    \textit{Repository} & \url{\af@repository}\\[1pt]
    \ifx\af@correspondence\@empty\else\textit{Correspondence} & \af@correspondence\\[1pt]\fi
    \textit{Disclosure} & \af@disclosure
  \end{tabular}\par
  \endgroup}

% Front matter: title page, then contents, then the body on a fresh page.
\newcommand{\affrontmatter}{%
  \maketitle
  \vspace*{\fill}
  \afpublicationblock
  \thispagestyle{plain}%
  \clearpage
  {\hypersetup{linkcolor=black}\fontsize{9}{10.6}\selectfont\tableofcontents}
  \clearpage}

% Unnumbered back-matter sections (amsart lists starred sections in the contents).
\newcommand{\afbacksection}[1]{\global\af@appendixpagefalse\section*{#1}}

\makeatother
% Notation local to this paper. Define only what materially simplifies the exposition.
\newcommand{\Z}{\mathbb{Z}}
\newcommand{\Q}{\mathbb{Q}}
\newcommand{\R}{\mathbb{R}}
\newcommand{\C}{\mathbb{C}}

% Special functions and arithmetic
\newcommand{\Li}{\operatorname{Li}}
\DeclareMathOperator{\cont}{cont}
\DeclareMathOperator{\real}{Re}                % real part
\DeclareMathOperator{\imag}{Im}                % imaginary part
\newcommand{\vg}[1]{v^{\mathrm{G}}_{#1}}      % Gauss valuation at the prime #1
\newcommand{\Zp}{\mathbb{Z}_{(p)}}            % integers localised at p
\newcommand{\ind}{\mathbf{1}}

% Objects of the construction
\newcommand{\Qt}{\widetilde{Q}}               % normalised determinant
\newcommand{\cR}{\mathcal{R}}                 % admissible rational functions
\newcommand{\cV}{\mathcal{V}}                 % local functional
\newcommand{\cT}{\mathcal{T}}                 % reflection tail
\newcommand{\cS}{\mathcal{S}}                 % slot multisets
\newcommand{\cJ}{\mathcal{J}}                 % jet bound
\newcommand{\cP}{\mathcal{P}}                 % pole bound
\newcommand{\Ex}{\mathcal{E}}                 % exceptional primes
\newcommand{\pP}{\mathfrak{p}}                % exceptional prime dividing 1-lambda
\newcommand{\pQ}{\mathfrak{q}}                % exceptional prime dividing the denominator of lambda
\newcommand{\rar}{\rho_{\mathrm{ar}}}         % arithmetic rate
\newcommand{\ran}{\rho_{\mathrm{an}}}         % analytic rate
\newcommand{\Cst}{C_{*}}                      % normalisation constant
\newcommand{\dd}{\,\mathrm{d}}
\newcommand{\eps}{\varepsilon}


\title[Enlarged Hankel determinants and four dilogarithm values]{Enlarged Hankel determinants and the irrationality of $\Li_2(2/11)$, $\Li_2(2/9)$, $\Li_2(-2/9)$ and $\Li_2(-2/7)$}
\hypersetup{pdftitle={Enlarged Hankel determinants and the irrationality of Li₂(2/11), Li₂(2/9), Li₂(−2/9) and Li₂(−2/7)}}
\author{AxiomForge}
\urladdr{https://github.com/axiomforge-math}
\subjclass[2020]{Primary 11J72; Secondary 33B30, 31A15}
\keywords{dilogarithm, irrationality, Hankel determinants, Heine's formula, logarithmic energy, prime number theorem, Landen's identity}
\date{26 November 2026}

\afidentifier{AF-0001}
\afversion{0.1.0}
\afdate{26 November 2026}
\afstatus{Preprint --- complete written argument, not yet independently verified. The finite inequalities behind the numerical rates have been reproduced and recomputed by a separately written implementation (an AI-internal check, not an independent human or external verification); no human review has taken place. Lean formalization in preparation: the main theorems are proved in Lean only conditionally on the three estimates.}
\afrepository{https://github.com/axiomforge-math/mathematics/tree/main/results/AF-0001-dilogarithm-irrationality}
\afcorrespondence{\href{mailto:axiomforge.math@gmail.com}{axiomforge.math@gmail.com}}
\afdisclosure{AI-assisted research. The original exploration, constructions, candidate proofs and certificate programs were produced with GPT-6.2 XHigh (web version); the reconstruction, the audit computations, the recomputation of the certificates, this manuscript and the Lean formalization with Claude Opus~5.5 (Anthropic); an adversarial review and the resulting revision with Claude Fable~5.2 (Anthropic); all under the direction of AxiomForge. No independent human review has taken place. A detailed AI-assisted research disclosure appears at the end.}

\begin{document}

\begin{abstract}
We present a candidate proof that the four numbers $\Li_2(2/11)$, $\Li_2(2/9)$, $\Li_2(-2/9)$ and $\Li_2(-2/7)$ are irrational. These values lie below the ranges covered by the published and the recent preprint literature that we could locate. The two Landen pairs $\{\Li_2(2/11),\Li_2(-2/9)\}$ and $\{\Li_2(2/9),\Li_2(-2/7)\}$ each contain an irrational number for elementary reasons. The argument uses the normalized Hankel determinants of Tang, in the form written up by Litsey, with a larger family of weights $D_{Bn}(x)^3/D_{Cn}(x)$ and dimensions $Hn$, the parameters $(H,B,C)$ depending on the value. Three estimates are derived for the resulting integer polynomials, with written proofs that have not been independently verified: a lower bound for their content, obtained prime by prime through a distribution of the functional over residue classes and integrated with the prime number theorem; an upper bound for their value at the dilogarithm, from Heine's formula and a comparison of logarithmic energies (on the real axis for the positive arguments and on a pair of rays for the negative ones); and nonvanishing at every fixed rational point for all large prime indices, by reduction to two explicit local Hankel determinants. The finite inequalities on which the numerical rates rest are checked in exact rational arithmetic. For $\Li_2(2/9)$ the two rates differ by only $0.019$ per squared dimension, and the exact data available for small $n$ lie above the upper bound and do not corroborate it. The proof is not formalized; we state precisely which steps have been checked and how.
\end{abstract}

\affrontmatter

\section{Introduction}\label{sec:introduction}

For $|z|\le 1$ let $\Li_2(z)=\sum_{k\ge1}z^k/k^2$ be the dilogarithm. This paper presents a proof of the following statement. The proof is a candidate: it has not been independently verified, and its status is described precisely in \cref{sec:intro-status}.

\begin{theorem}\label{thm:main}
The numbers $\Li_2(2/11)$, $\Li_2(2/9)$, $\Li_2(-2/9)$ and $\Li_2(-2/7)$ are irrational.
\end{theorem}

\subsection{Context}\label{sec:intro-context}
Irrationality of the dilogarithm at rational points $r/s$ is classical for $|s|$ large compared with $|r|$. Hata proved irrationality and linear independence results for values $\Li_2(1/q)$ \cite{Hata1990,Hata1993}, and the permutation group method of Rhin and Viola \cite{RhinViola2005,RhinViola2019} and the work of Viola and Zudilin \cite{ViolaZudilin2018} extended these results to other rational arguments, with explicit thresholds on the denominator. For numerator $2$, the method of \cite{RhinViola2005} gives the irrationality of $\Li_2(2/s)$ for $s\ge51$, as stated in \cite[Section~1]{TanZhang2026} and \cite[Section~1.2]{Litsey2026Li2}; we have not checked this threshold in the original, and we know of no refereed result for the denominators $7$, $9$ and $11$ considered here.

In 2026 a different construction was introduced in Tang's Lean development for $\Li_2(1/2)$ \cite{TangLi2Half}: normalized Hankel determinants of an affine moment functional, an instance of the determinantal approach to irrationality \cite{Zudilin2017}. Tang's development is Lean code without a paper; we know its construction through the written account in \cite[Sections~2 and~9]{Litsey2026Li2}. The method was extended, in unrefereed and largely AI-assisted work, by Liu \cite{LiuLi2Rational}, whose Lean development states the irrationality of $\Li_2(2/k)$ for $k\ge13$ and of $\Li_2(-2/k)$ for $k\ge16$, and it is used for small numerators by Tan and Zhang \cite{TanZhang2026}, whose Table~1 gives the ranges $q\ge12$ and $q\ge16$ for $\Li_2(2/q)$ and $\Li_2(-2/q)$. Prellberg \cite{Prellberg2026} treated the arguments $-1/4$, $1/5$ and $-1/3$ by the Rhin--Viola and Viola--Zudilin constructions. Litsey's preprint \cite{Litsey2026Li2} proves irrationality measures with Tang's determinants; its Corollary~1.3 covers $\Li_2(2/b)$ for every odd $b\ge13$ and $\Li_2(-2/b)$ for every odd $b\ge11$, its Section~1.1 states that its main theorem gives no measure for $\Li_2(2/b)$ with $b\in\{5,7,9,11\}$ or for $\Li_2(-2/b)$ with $b\in\{5,7,9\}$, and its Remark~8.1 records that the margin of that construction is at most $0$ at $2/11$ and at $-2/9$. We read these five sources on 26 November 2026; none of them is refereed. We know of no source, refereed or not, that proves the irrationality of any of the four numbers in \cref{thm:main}; this is a statement about what we found, not a proof of absence.

The four numbers come in two pairs linked by Landen's identity, and the following weaker statement is elementary.

\begin{proposition}\label{prop:landen}
At least one of $\Li_2(2/11)$, $\Li_2(-2/9)$ is irrational, and at least one of $\Li_2(2/9)$, $\Li_2(-2/7)$ is irrational.
\end{proposition}

\begin{proof}
For real $z<1/2$ Landen's identity reads
\begin{equation}\label{eq:landen}
\Li_2(z)+\Li_2\Bigl(\frac{z}{z-1}\Bigr)=-\tfrac12\log^2(1-z).
\end{equation}
Indeed both sides vanish at $z=0$, and with $w=z/(z-1)$ one has $1-w=(1-z)^{-1}$ and $w'=-(z-1)^{-2}$, so the derivative of the left side is
$-z^{-1}\log(1-z)-\log(1-z)\,(z(z-1))^{-1}=\log(1-z)/(1-z)$, which is the derivative of the right side. For $z=2/11$ and $z=2/9$ we have $z/(z-1)=-2/9$ and $-2/7$, and \eqref{eq:landen} gives
\[
\Li_2\bigl(\tfrac2{11}\bigr)+\Li_2\bigl(-\tfrac29\bigr)=-\tfrac12\log^2\tfrac9{11},\qquad
\Li_2\bigl(\tfrac29\bigr)+\Li_2\bigl(-\tfrac27\bigr)=-\tfrac12\log^2\tfrac79 .
\]
By the Hermite--Lindemann theorem $\log r$ is transcendental for every rational $r>0$, $r\ne1$; hence so is $\log^2 r$, and the right-hand sides are irrational.
\end{proof}

\Cref{prop:landen} cannot be sharpened by this route, and \cref{thm:main} is a statement about the individual values. The four arguments are exactly the two Landen pairs $(z,z/(z-1))$ for $z\in\{2/11,2/9\}$.

\subsection{Relation to earlier work}\label{sec:intro-sources}
The method of this paper is that of Tang and Litsey, applied to a larger family of weights, and we make no claim of novelty for its techniques. The following ingredients are taken from \cite{Litsey2026Li2}, where they are attributed in part to Tang \cite{TangLi2Half}; section numbers are those of \cite{Litsey2026Li2}.
\begin{itemize}
\item The functional $\Lambda_X$, the Hankel determinants, their normalization and Heine's formula (Section~2). Tang's family is the case $(H,B,C)=(2,1,4)$ of \cref{sec:intro-construction}.
\item The arithmetic at the primes dividing $b$ and $b-2\epsilon$ and at the other primes: the bound through integer-valued polynomials with the parameter $f$ (Lemma~4.2 and Theorem~4.1), the reflection formula and the pole slots at a prime dividing $b$ (Lemma~9.9 and Proposition~9.8), the determinant lemmas (Lemmas~9.4 and~9.5), the residue classes modulo $p$ with the local exponents $3a_\sigma-c_\sigma-1+\ind_{\sigma=0}$ and the resulting bounds at the primes $p\le n$ and $n<p\le4n$ (Propositions~9.11, 9.12 and~9.15), and the integration of the resulting profile with the prime number theorem (Proposition~4.3 and Sections~9.10--9.11). Our \cref{sec:exceptional,sec:ordinary,sec:profile} adapt these arguments to general $(H,B,C)$; the jet bound and the pole bound of \cref{prop:jet,prop:pole} are our formulation, for general parameters, of the method of Sections~9.8 and~9.9, and the proof of the local integrality through finite geometric averages (\cref{prop:distribution,lem:local-integrality}) differs from the proof given there.
\item Nonvanishing through residue classes modulo a prime and local Hankel determinants of a local functional (Section~5, where they are called type constants). There the prime runs through a band depending on $n$; here we take the index $n$ itself prime, which leaves only two local determinants (\cref{sec:nonvanishing}).
\item The analytic bound: Heine's formula for a functional with jets, with at most two variables at a site (Lemma~6.2), the comparison with the energy of a measure, and the bound with a level $\theta$ and the penalty $2\int(\cdot-\theta)_+$ (Theorem~6.1); for negative arguments, the replacement of the sum over the integers by a contour integral along a bent contour (Section~6.3). Our \cref{prop:positive-bound} has the form of Theorem~6.1 for the lattice; in \cref{sec:contour} the contour consists of two rays from $\frac12$, and the energy inequality (\cref{lem:energy-comparison}) is proved differently.
\end{itemize}
What is specific to this paper is the choice of the parameters $(H,B,C)$ of \cref{tab:parameters}, in particular the family without zeros for $2/9$, the statements for general parameters, the reduction at prime indices, the comparison measures and the finite computations of \cref{app:certificates}. We have compared our statements with \cite{Litsey2026Li2} only at the level of this list, not line by line, and we have not read Tang's or Liu's Lean code.

\subsection{The construction}\label{sec:intro-construction}
Fix one of the four arguments $\lambda$, write $\xi=\Li_2(\lambda)$, and let $(H,B,C)$ be the triple of integers attached to $\lambda$ in \cref{tab:parameters}. For $n\ge1$ put
\[
h=Hn,\qquad A=Bn,\qquad K=Cn,\qquad D_m(x)=\prod_{j=1}^m(x+j),\qquad W_n(x)=\frac{D_A(x)^3}{D_K(x)} .
\]
The weight $W_n$ has double zeros at $-1,\dots,-A$ and simple poles at $-A-1,\dots,-K$. Let $\Lambda_X$ be the $\Q[X]$-valued linear functional on rational functions with at most simple poles at negative integers defined in \cref{def:functional}; its value at $X=\xi$ is
\[
\Lambda_\xi(R)=\sum_{x\ge1}\lambda^x\bigl(xR(x)\bigr)' .
\]
The object of study is the polynomial
\begin{equation}\label{eq:intro-Q}
\Qt_n(X)=\frac{S_n^{\,h}}{F_n}\det_{0\le i,j<h}\Lambda_X\bigl(W_n(x)\,x^{i+j}\bigr),\qquad S_n=\frac{K!}{(A!)^3},\quad F_n=\prod_{i=0}^{h-1}(i!)^2,
\end{equation}
which has rational coefficients and degree at most $h$. Tang's family, as described in \cite[Section~2]{Litsey2026Li2}, is the case $(H,B,C)=(2,1,4)$. The families used here are larger: the dimension $h$ and the pole range $K$ grow faster than the numerator length $A$, and for $\lambda=2/9$ the weight has no zeros at all.

\begin{table}[ht]
\centering
\caption{Parameters. Here $\lambda=2\epsilon/b$ with $\epsilon=\pm1$; $\pP$ is the prime dividing the numerator $b-2\epsilon$ of $1-\lambda$, with $f=v_{\pP}(b-2\epsilon)$, and $\pQ$ is the prime dividing $b$, with $e=v_{\pQ}(b)$. The normalized parameters are $a=B/H$, $\kappa=C/H$ and $d=\kappa-3a$.}
\label{tab:parameters}
\begin{tabular}{@{}lccccccccc@{}}
\toprule
$\lambda$ & $(H,B,C)$ & $a$ & $\kappa$ & $d$ & $\pP$ & $f$ & $\pQ$ & $e$ & $t=\log(b/2)$\\
\midrule
$2/11$ & $(4,1,6)$ & $1/4$ & $3/2$ & $3/4$ & $3$ & $2$ & $11$ & $1$ & $\log(11/2)$\\
$2/9$ & $(8,0,9)$ & $0$ & $9/8$ & $9/8$ & $7$ & $1$ & $3$ & $2$ & $\log(9/2)$\\
$-2/9$ & $(4,1,7)$ & $1/4$ & $7/4$ & $1$ & $11$ & $1$ & $3$ & $2$ & $\log(9/2)$\\
$-2/7$ & $(5,1,8)$ & $1/5$ & $8/5$ & $1$ & $3$ & $2$ & $7$ & $1$ & $\log(7/2)$\\
\bottomrule
\end{tabular}
\end{table}

\subsection{The argument in one page}\label{sec:one-page}
For a nonzero $F\in\Q[X]$ let $\cont(F)>0$ be the rational number for which $F/\cont(F)$ is a primitive polynomial in $\Z[X]$ (\cref{sec:content}). The proof rests on the following elementary criterion.

\begin{lemma}[irrationality criterion]\label{lem:criterion}
Let $x\in\R$, let $(F_k)_{k\ge1}$ be polynomials in $\Q[X]$ with $\deg F_k\le h_k$, where $h_k\to\infty$, and let $\rho>\rho'$ be real numbers. Assume that for every $r\in\Q$ one has $F_k(r)\ne0$ for all large $k$ (in particular $F_k\ne0$ for large $k$), and that for every $\eta>0$ and all large $k$
\[
\log\cont(F_k)\ge(\rho-\eta)h_k^2\qquad\text{and}\qquad\log|F_k(x)|\le(\rho'+\eta)h_k^2 .
\]
Then $x$ is irrational.
\end{lemma}

\begin{proof}
Suppose $x=u/v$ with $u\in\Z$, $v\in\Z_{\ge1}$, and take $\eta=\frac13(\rho-\rho')$. For large $k$ the polynomial $P_k=F_k/\cont(F_k)\in\Z[X]$ satisfies $P_k(u/v)\ne0$, so $v^{h_k}P_k(u/v)$ is a nonzero integer and $0\le h_k\log v+\log|F_k(x)|-\log\cont(F_k)\le h_k\log v-\eta h_k^2$, which is negative for large $k$.
\end{proof}

We apply \cref{lem:criterion} to $x=\xi$ and $F_k=\Qt_{p_k}$, where $p_k$ is the $k$th prime and $h_k=Hp_k$. Its three hypotheses are the following theorems, in which $\rar(\lambda)$ and $\ran(\lambda)$ are the constants of \cref{tab:rates}.

\begin{theorem}[arithmetic estimate]\label{thm:arith}
$\displaystyle\liminf_{n\to\infty,\ \Qt_n\ne0}\ h^{-2}\log\cont(\Qt_n)\ \ge\ \rar(\lambda)$.
\end{theorem}

\begin{theorem}[archimedean estimate]\label{thm:arch}
$\displaystyle\limsup_{n\to\infty}\ h^{-2}\log|\Qt_n(\xi)|\ \le\ \ran(\lambda)$, with the convention $\log0=-\infty$.
\end{theorem}

\begin{theorem}[nonvanishing at prime indices]\label{thm:nonvanishing}
For every $r\in\Q$ there is $p_0(r)$ such that $\Qt_p(r)\ne0$ for every prime $p\ge p_0(r)$.
\end{theorem}

\begin{table}[ht]
\centering
\caption{Thresholds and computed bounds, normalized by $h^2$. The thresholds $\rar$ and $\ran$ are the rational numbers used in \cref{thm:arith,thm:arch}. What is checked in exact arithmetic is a set of finite inequalities: that the computed lower bound for $I_\delta+\eps_{\mathrm{exc}}$ (\cref{tab:profile}) exceeds $\rar$, and that the computed upper bound for the comparison-measure functional (\cref{tab:analytic}) is below $\ran$; that these quantities bound the asymptotic rates is the content of the candidate proofs. The column $\rar-\ran$ is the margin that enters \cref{lem:criterion}; ``computed'' is the difference of the two computed bounds; ``recomputed'' is the margin obtained for the same comparison measures by the AI-internal recomputation, with exact profile integrals and finer cell enclosures (\cref{sec:margins}). Computed lower bounds are rounded down, upper bounds up, and margins to the nearest digit. For $\lambda=2/11$ the certificates work in the scale $n^2=h^2/16$ (\cref{rem:211-scale}). For $\lambda=2/9$ the margin is thin, and the exact data for $n\le11$ lie above $\ran$: this case has no numerical corroboration (\cref{sec:margins}).}
\label{tab:rates}
\footnotesize
\setlength{\tabcolsep}{3.5pt}
\begin{tabular}{@{}lccccccc@{}}
\toprule
 & & lower bound & & upper bound & & \multicolumn{2}{c}{margin of the bounds}\\
\cmidrule(l){7-8}
$\lambda$ & $\rar$ & computed & $\ran$ & computed & $\rar-\ran$ & computed & recomputed\\
\midrule
$2/11$ & $-2169/8000$ & $-0.2710651$ & $-49/160$ & $-0.3071494$ & $0.035125$ & $0.036084$ & $0.0678$\\
$2/9$ & $-133/100$ & $-1.3245275$ & $-134/100$ & $-1.3436316$ & $0.010000$ & $0.019104$ & $0.019379$\\
$-2/9$ & $-2/5$ & $-0.3830882$ & $-49/100$ & $-0.4954089$ & $0.090000$ & $0.112321$ & $0.13978$\\
$-2/7$ & $-521323/10^6$ & $-0.5213227$ & $-526078/10^6$ & $-0.5260781$ & $0.004755$ & $0.0047556$ & $0.02969$\\
\bottomrule
\end{tabular}
\end{table}

Since $\rar>\ran$ in each case, \cref{lem:criterion,thm:arith,thm:arch,thm:nonvanishing} prove \cref{thm:main} (\cref{sec:proof-main}). The three estimates are independent of one another, and each is obtained as follows. Both $\log\cont(\Qt_n)$ and $\log|\Qt_n(\xi)|$ are of order $h^2$ only because of the factor $S_n^h/F_n$ in \eqref{eq:intro-Q}: its logarithm is $(d-1)h^2\log h+\Cst h^2+O(h\log h)$ (\cref{lem:stirling}), which cancels the term $(1-d)h^2\log h$ of $\log|Q_n(\xi)|$ in the archimedean estimate, and its $p$-adic valuation is the term $h\mathsf N(p/h)$ of the arithmetic profile (\cref{def:profile}).

\emph{Arithmetic (\cref{sec:exceptional,sec:ordinary,sec:profile}).} In the binomial basis, $\Qt_n$ is a Gram determinant whose entries have integral residues (\cref{lem:normalization,lem:integral-valued}), and $\log\cont(\Qt_n)=\sum_p\vg p(\Qt_n)\log p$ is bounded prime by prime: at the three exceptional primes of \cref{tab:parameters} directly (\cref{sec:exceptional}); at an ordinary prime $p$ with $\sqrt K<p\le K$ by distributing $\Lambda_X$ over the residue classes modulo $p$ and expanding the determinant by Cauchy--Binet (\cref{prop:jet,prop:pole}). For $p\ge\delta h$, with a fixed small cut-off $\delta>0$, the resulting bounds are $h$ times a function of $p/h$, up to a bounded error (\cref{prop:profile-uniform}); the prime number theorem integrates them (\cref{thm:arith-general}), the ordinary primes below $\delta h$ are bounded crudely at the cost $-2\delta$ in the rate, and the integral is bounded below by an exact computation (\cref{prop:profile-certificate}).

\emph{Archimedean (\cref{sec:energy,sec:contour}).} Heine's formula writes $Q_n(\xi)$ as an average of a squared Vandermonde determinant against a product weight, over $h$-tuples of positive integers for $\lambda>0$ and, after a residue identity that keeps the alternating sign, over $h$-tuples of points on two rays for $\lambda<0$ (\cref{lem:contour}); a regularized logarithmic-energy inequality against an explicit comparison measure bounds this average (\cref{prop:positive-bound,prop:negative-bound,tab:analytic}). For $\lambda>0$ the functional involves a derivative, and a positive integer can carry at most two of the $h$ variables (\cref{lem:multiplicity}); this occupancy constraint is what limits the gain, and it enters the bound as the penalty $2\int(\mathsf W-\theta)_+$ of \cref{prop:positive-bound}.

\emph{Nonvanishing (\cref{sec:nonvanishing}).} At a prime index $n=p$, in a basis adapted to the Chinese remainder theorem modulo $p$, the scaled Gram matrix at a rational point is block triangular modulo $p$, with diagonal blocks equal up to units to two fixed $H\times H$ Hankel matrices whose determinants are nonzero rational numbers (\cref{thm:nonvanishing-general,tab:local-dets}).

Besides elementary facts (Legendre's and Kummer's formulas, Stirling's formula, Newton interpolation, the Cauchy--Binet formula, Schur's product theorem, Fubini's theorem), the proof of \cref{thm:main} uses two external theorems, the prime number theorem in the form $\sum_{p\le y}\log p\sim y$ and the residue theorem, and three kinds of finite computations in exact rational arithmetic: four lower bounds for profile integrals, eight local determinants and four upper bounds for comparison-measure functionals (\cref{app:certificates}). The Hermite--Lindemann theorem is used only in \cref{prop:landen}.

\emph{The case $\lambda=2/9$.} The four cases are not equally robust. For $2/11$, $-2/9$ and $-2/7$ the exact values of $h^{-2}\log|\Qt_n(\xi)|$ computed for small $n$ lie below $\ran$, and the recomputed margins are between $0.03$ and $0.14$. For $2/9$ the computed bounds are $-1.3245$ and $-1.3436$, against rates of size $1.3$: the margin is $0.019$ per $h^2$, and $0.010$ between the thresholds. The finite data do not corroborate the upper bound: at $n=11$ the value $h^{-2}\log|\Qt_n(\xi)|=-1.2886$ is still above $\ran$ and above $\rar$ (\cref{tab:numerics}), and the comparison measure is nearly optimal for the bound of \cref{prop:positive-bound}, so the certificate cannot be improved much (\cref{sec:margins}). The case $2/9$ therefore rests on the written proofs alone, and an unnoticed loss of $0.02\,h^2$ in any step of \cref{sec:exceptional,sec:ordinary,sec:profile,sec:energy} would invalidate it.

\subsection{Status of the argument}\label{sec:intro-status}
The proof is a candidate. The choice of the families, the arguments and the certificates were first produced by an AI system, as three separate candidate manuscripts, on the basis of the method described in \cref{sec:intro-sources}; this paper reconstructs them as one argument, corrects them where necessary, and states each step with its hypotheses. \Cref{tab:status} records the status of each step in three categories:
\begin{itemize}
\item \emph{proved}: an elementary argument, given in full, of a kind we regard as standard;
\item \emph{certified}: a finite inequality or identity checked by a computation in exact rational arithmetic, reproduced by the original programs and confirmed by an AI-internal recomputation from the definitions, which is not an independent or human check; the word applies to the finite statement only, not to the asymptotic rate deduced from it;
\item \emph{candidate}: a complete proof is given here, but it has not been independently verified.
\end{itemize}
The candidate steps were checked only in the following limited sense. A separate exact implementation of the construction, written from the definitions, confirms for all four values and for $n$ up to $11$--$19$ that $Q_n\ne0$, that the series identity of \cref{lem:series} holds to forty digits, and that every exact valuation bound of \cref{sec:exceptional,sec:ordinary} holds at every prime (\cref{app:numerics}). This audit found one error in the source argument: there, the pole bound fails by exactly one unit, because the pole $l=p$ of the zero residue class gains only $+1$ rather than $+2$; \cref{prop:pole,rem:zero-class} state the corrected bound, and the jet bound is unaffected. All checks were carried out with AI systems; none of them is an independent verification.

\begin{table}[ht]
\centering
\caption{Status of the steps of the proof.}
\label{tab:status}
\small
\begin{tabular}{@{}>{\raggedright\arraybackslash}p{0.44\textwidth}>{\raggedright\arraybackslash}p{0.2\textwidth}>{\raggedright\arraybackslash}p{0.28\textwidth}@{}}
\toprule
Step & Where & Status\\
\midrule
Landen pairs & \cref{prop:landen} & proved\\
Irrationality criterion & \cref{lem:criterion} & proved; formalized\\
Series identity, binomial normalization, degree bound & \cref{lem:series,lem:normalization} & proved\\
Exceptional primes and uniform bounds & \cref{sec:exceptional} & candidate; finite bounds checked on exact data\\
Distribution over residue classes, local integrality, jet and pole bounds & \cref{sec:ordinary} & candidate; jet and corrected pole bounds checked on exact data\\
Uniform profile and prime number theorem summation & \cref{prop:profile-uniform,thm:arith-general} & candidate; the comparison of the discrete and limiting profiles checked on exact class counts\\
Lower bounds for the profile integrals & \cref{prop:profile-certificate} & certified\\
Nonvanishing reduction at prime indices & \cref{thm:nonvanishing-general} & candidate\\
Local determinants & \cref{prop:local-dets} & certified\\
Heine formula, multiplicity bound, energy comparison & \cref{sec:energy} & candidate\\
Contour representation and bound for $\lambda<0$ & \cref{sec:contour} & candidate\\
Upper bounds for the comparison-measure functionals & \cref{tab:analytic} & certified\\
Final deduction & \cref{sec:proof-main} & proved; formalized conditionally\\
\bottomrule
\end{tabular}
\end{table}

No independent human review of the general arguments and no independent external check of the certificates have taken place. The Lean development formalizes the definitions and the general irrationality criterion, and proves the four main theorems only conditionally, from hypotheses that bundle \cref{thm:arith,thm:arch,thm:nonvanishing}; these three estimates are not formalized. The margins of \cref{tab:rates} are small for $\lambda=2/9$ and $\lambda=-2/7$. For $-2/7$ this reflects generous error bounds in the certificate, and the recomputation gives a margin of $0.02969$; for $2/9$ the margin is intrinsically small and is not supported by the finite data (\cref{sec:one-page,sec:margins}). The manuscript was revised after an adversarial review by a further AI system, which found no false step but several gaps and imprecisions; that review is an AI review and not an independent one. \Cref{sec:open} lists the points that we regard as open or as deserving particular scrutiny.

\subsection{Organization, notation and conventions}\label{sec:conventions}
\Cref{sec:functional} defines the functional and the determinants; \cref{sec:exceptional,sec:ordinary,sec:profile} prove \cref{thm:arith}; \cref{sec:nonvanishing} proves \cref{thm:nonvanishing}; \cref{sec:energy,sec:contour} prove the comparison bounds behind \cref{thm:arch}; \cref{sec:proof} assembles the proof and discusses margins and open points. \Cref{app:certificates} describes the certificates, \cref{app:technical} collects technical lemmas, and \cref{app:numerics} reports numerical evidence for small $n$.

$v_p$ is the $p$-adic valuation with $v_p(p)=1$, and $\Zp$ is the ring of rational numbers with denominator prime to $p$. For $F\in\Q[X]$, $\vg p(F)$ is the minimum of the valuations of its coefficients, with $\vg p(0)=+\infty$. We write $(y)_+=\max(y,0)$, $\ind_{\mathsf{S}}$ for the indicator of a statement $\mathsf S$, $b_i(x)=\binom xi$, and $\lfloor\cdot\rfloor$ for the integer part; logarithms are natural. Throughout, $\lambda$ is one of the four arguments, and all quantities attached to the construction depend on $\lambda$ and $n$ even when this is not indicated. \Cref{tab:notation} lists the principal symbols and where they are defined. A few letters have, besides the global meaning of \cref{tab:notation}, a local meaning that is stated where it is used: $q$ (a generic rational argument in \cref{sec:moments}, the polynomial part of a rational function, and the integers $q_s$, $q_{ij}$), $c$ (residues $c_l$, class counts $c_\sigma$, the constant $c_R$, and the ray direction of \cref{sec:contour}), $B$ and $D$ (the parameter $B$ and the products $D_m$, but also the denominators $B(u)$, $D(u)$ of \cref{lem:local-integrality}, the functional $D_v$ and index sets $D$), $T$ (the averages $T_{M,v}$, a triangular matrix in \cref{lem:normalization}, and the decay rate of \cref{sec:contour}), $N$ (the integer $2h+3A$, and a truncation index in \cref{lem:contour,app:certificates}), $t$ ($\log(b/2)$, and an integer or a price in \cref{lem:integral-valued,lem:weak-duality}), $e$ (the exponent of \cref{tab:parameters}, the basis $e_{\sigma,i}$ and the exponent $e_0$ of \cref{sec:nonvanishing}), and the constants $C_0$, $C_1$, $c_1$, $Y$, which are local to the statement in which they occur.

\begin{table}[ht]
\centering
\caption{Notation.}
\small
\label{tab:notation}
\small
\begin{tabular}{@{}>{\raggedright\arraybackslash}p{0.27\textwidth}>{\raggedright\arraybackslash}p{0.44\textwidth}>{\raggedright\arraybackslash}p{0.21\textwidth}@{}}
\toprule
Symbol & Meaning & Defined in\\
\midrule
$\lambda=2\epsilon/b$, $\xi$, $t$ & argument, $\xi=\Li_2(\lambda)$, $t=\log(b/2)$ & \cref{tab:parameters}\\
$H,B,C$; $h,A,K$ & parameters; $h=Hn$, $A=Bn$, $K=Cn$ & \cref{sec:intro-construction}\\
$a,\kappa,d$ & $B/H$, $C/H$, $\kappa-3a$ & \cref{tab:parameters}\\
$\pP,f$; $\pQ,e$; $\Ex$ & exceptional primes and exponents; $\Ex=\{2,\pP,\pQ\}$ & \cref{tab:parameters}\\
$D_m$, $W_n$, $S_n$, $F_n$ & rising product, weight, normalizing factors & \eqref{eq:intro-Q}\\
$M_k$, $\tau_l$, $G_q$, $G$ & geometric moments, partial sums of $\Li_2$, geometric sums & \cref{sec:moments}\\
$\cR$, $\Lambda_X$ & admissible rational functions, the functional & \cref{def:functional}\\
$Q_n$, $\Qt_n$, $R_{ij}$ & determinant, normalized determinant, binomial entries & \cref{sec:determinants}\\
$N$, $L_p$ & $2h+3A$, $\lfloor\log N/\log p\rfloor$ & \eqref{eq:N-L}\\
$\cont$, $\vg p$ & content, Gauss valuation & \cref{sec:content}\\
$\lambda_p$; $a_\sigma,c_\sigma,w_\sigma$ & $\lambda^p$; class counts modulo an ordinary prime & \cref{sec:ordinary}\\
$T_{M,v}$, $T_{M,v,Y}$, $Y_\mu$, $\Theta_\sigma$ & finite averages, regularization, class terms & \cref{sec:finite-averages,sec:distribution}\\
$\cJ_p$, $\cP_p$ & jet and pole bounds & \cref{prop:jet,prop:pole}\\
$J,R,\mathsf N,\mathsf B$, $I_\delta$ & profile and its integral & \cref{def:profile,thm:arith-general}\\
$\eps_{\mathrm{exc}}$ & exceptional rate & \cref{cor:exceptional}\\
$\cV$, $w_0,w_1$, $\mathsf H_0,\mathsf H_1$ & local functional, weights, Hankel matrices & \cref{sec:local-functional}\\
$U_\sigma$, $E(\sigma)$ & potential and energy of a comparison measure & \cref{sec:energy-inequality}\\
$\Cst$ & normalization constant & \cref{lem:stirling}\\
$V$, $\mathsf W$; $V_-$, $c$, $T$, $k$, $\Gamma$ & external fields; ray direction, decay, kernel, contour & \cref{sec:positive-bound,sec:contour}\\
$\rar,\ran$ & thresholds for the two rates & \cref{tab:rates}\\
\bottomrule
\end{tabular}
\end{table}

\section{The affine functional and the Hankel family}\label{sec:functional}

This section defines the functional $\Lambda_X$ and the determinants $Q_n$, $\Qt_n$, and proves the three facts about them that the rest of the paper uses: the series identity at $X=\xi$ (\cref{lem:series}, used in \cref{sec:energy,sec:contour}), the expression of $\Qt_n$ as a Gram determinant in the binomial basis (\cref{lem:normalization}), and the integrality of its entries (\cref{lem:integral-valued}, used in \cref{sec:exceptional}). All statements here are elementary.

\subsection{Moments and the functional}\label{sec:moments}
For a rational number $q$ with $|q|<1$ and an integer $k\ge0$ put
\[
M_k(q)=\sum_{u\ge1}q^uu^k,\qquad \tau_l(q)=\sum_{u=1}^{l}\frac{q^u}{u^2}\quad(l\ge0),
\]
so that $\tau_0=0$. Since $\sum_{u\ge0}\binom uj q^u=q^j(1-q)^{-j-1}$ and $x^k=\sum_{j}{k\brace j}j!\binom xj$, where ${k\brace j}$ are the Stirling numbers of the second kind,
\begin{equation}\label{eq:moments}
M_0(q)=\frac{q}{1-q},\qquad M_k(q)=\sum_{j=1}^k{k\brace j}\,j!\,\frac{q^j}{(1-q)^{j+1}}\quad(k\ge1).
\end{equation}
For a polynomial $g$ we write $G_q(g)=\sum_{u\ge1}q^ug(u)$; by \eqref{eq:moments} it is a rational number. When $q=\lambda$ we drop the subscript and write $G=G_\lambda$.

Let $\cR$ be the $\Q$-vector space of rational functions $R\in\Q(x)$ whose poles are simple and lie in $\{-1,-2,\dots\}$. Every $R\in\cR$ has a unique partial fraction decomposition $R=q+\sum_{l\ge1}c_l\,(x+l)^{-1}$ with $q\in\Q[x]$ and finitely many nonzero $c_l\in\Q$. We call $q$ the polynomial part of $R$ and $c_l$ its residue at $-l$.

\begin{definition}\label{def:functional}
The functional $\Lambda_X\colon\cR\to\Q[X]$ is the $\Q$-linear map with
\[
\Lambda_X(x^k)=(k+1)M_k(\lambda)\quad(k\ge0),\qquad \Lambda_X\Bigl(\frac1{x+l}\Bigr)=l\lambda^{-l}\bigl(X-\tau_l(\lambda)\bigr)\quad(l\ge1).
\]
\end{definition}

Equivalently, for $R=q+\sum_lc_l(x+l)^{-1}$,
\begin{equation}\label{eq:functional-explicit}
\Lambda_X(R)=G\bigl((xq)'\bigr)+\sum_{l}l\,c_l\,\lambda^{-l}\bigl(X-\tau_l(\lambda)\bigr),
\end{equation}
since $(x\cdot x^k)'=(k+1)x^k$. The value of $\Lambda_X(R)$ is an affine polynomial in $X$.

\begin{lemma}[series identity]\label{lem:series}
For every $R\in\cR$ the series $\sum_{x\ge1}\lambda^x(xR(x))'$ converges absolutely, and
\[
\Lambda_\xi(R)=\sum_{x\ge1}\lambda^x\bigl(xR(x)\bigr)',\qquad \xi=\Li_2(\lambda).
\]
\end{lemma}

\begin{proof}
The derivative $(xR(x))'$ is a rational function without poles at the positive integers, so it is bounded by a polynomial in $x$ there, while $|\lambda|^x$ decays geometrically; this gives absolute convergence. By linearity it suffices to treat $R=x^k$ and $R=(x+l)^{-1}$. In the first case $(xR)'=(k+1)x^k$ and the identity is the definition of $M_k$. In the second case $(x/(x+l))'=l/(x+l)^2$, and the substitution $m=x+l$ gives
\[
\sum_{x\ge1}\lambda^x\frac{l}{(x+l)^2}=l\lambda^{-l}\sum_{m>l}\frac{\lambda^m}{m^2}=l\lambda^{-l}\bigl(\xi-\tau_l(\lambda)\bigr). \qedhere
\]
\end{proof}

\subsection{The determinants}\label{sec:determinants}
With the notation of \cref{sec:intro-construction} define
\begin{equation}\label{eq:Qn}
Q_n(X)=\det_{0\le i,j<h}\Lambda_X\bigl(W_nx^{i+j}\bigr),\qquad \Qt_n(X)=\frac{S_n^{\,h}}{F_n}\,Q_n(X).
\end{equation}
Since $W_nx^s\in\cR$ for every $s\ge0$, the entries are affine polynomials in $X$, and expanding the determinant over permutations shows $\deg Q_n\le h$.

\begin{lemma}[binomial normalization]\label{lem:normalization}
Put $R_{ij}=S_nW_nb_ib_j$ for $0\le i,j<h$. Then the polynomial $\Qt_n$ of \eqref{eq:Qn} equals $\det_{0\le i,j<h}\Lambda_X(R_{ij})$.
\end{lemma}

\begin{proof}
Write $b_i=\sum_{j<h}T_{ij}x^j$. The matrix $T$ is lower triangular with $T_{ii}=1/i!$, so $\det T=\prod_{i<h}(i!)^{-1}$. By bilinearity $[\Lambda_X(S_nW_nb_ib_j)]=S_n\,T\,[\Lambda_X(W_nx^{i+j})]\,T^{\mathsf T}$, whose determinant is $S_n^h(\det T)^2Q_n=\Qt_n$.
\end{proof}

The binomial form of the Gram matrix is the one in which the arithmetic is transparent. Put
\begin{equation}\label{eq:N-L}
N=2h+3A,\qquad L_p=\lfloor\log N/\log p\rfloor\quad(p\text{ prime}).
\end{equation}
The numerator of $R_{ij}$ has degree $3A+i+j<N$, and $p^{L_p}\le N<p^{L_p+1}$.

\begin{lemma}[integrality of the binomial entries]\label{lem:integral-valued}
Let $0\le i,j<h$ and write $R_{ij}=q_{ij}+\sum_{l=1}^Kc^{(ij)}_l(x+l)^{-1}$.
\begin{enumerate}[(a)]
\item $c^{(ij)}_l=\gamma_l\,b_i(-l)\,b_j(-l)$ with $\gamma_l=(-1)^{l-1}K\binom{K-1}{l-1}\binom{A-l}{A}^3\in\Z$, and $\gamma_l=0$ for $l\le A$.
\item $q_{ij}=\sum_{K\le k\le N}\beta_k\binom{x}{k-K}\binom kK^{-1}$ with $\beta_k\in\Z$; in particular $\deg q_{ij}\le i+j+3A-K$, and $q_{ij}=0$ if $i+j+3A<K$.
\item For every prime $p$, every $t\in\Z$ and every $s\ge0$, the coefficient of $y^s$ in $q_{ij}(t+y)$ has valuation at least $-(s+1)L_p$, and the coefficient of $y^s$ in $b_i(t+y)$ has valuation at least $-sL_p$.
\end{enumerate}
\end{lemma}

\begin{proof}
Since $D_A(x)=A!\binom{x+A}{A}$ and $D_K(x)=K!\binom{x+K}{K}$, we have $R_{ij}=K!\,\beta(x)/D_K(x)$ with $\beta=\binom{x+A}A^3b_ib_j$, an integer-valued polynomial of degree $3A+i+j$.

(a) The residue of $K!\beta/D_K$ at $-l$ is $K!\beta(-l)/\prod_{m\ne l}(m-l)=(-1)^{l-1}K\binom{K-1}{l-1}\beta(-l)$, and $\binom{-l+A}{A}$ vanishes for $1\le l\le A$.

(b) Expand $\beta=\sum_k\beta_k\binom{x+K}{k}$ with $\beta_k\in\Z$ (an integer-valued polynomial has integral coefficients in the basis $\binom{x+K}k$). For $k<K$ the quotient $\binom{x+K}k/\binom{x+K}K$ is a proper rational function. For $k\ge K$ one checks $\binom{x+K}k/\binom{x+K}K=\binom{x}{k-K}/\binom kK$. Hence the polynomial part of $K!\beta/D_K$ is as stated.

(c) By Kummer's theorem $v_p\binom kK$ is the number of carries in the addition of $K$ and $k-K$ in base $p$, which is at most $\lfloor\log k/\log p\rfloor\le L_p$ for $k\le N$. Hence $p^{L_p}q_{ij}$ takes values in $\Zp$ at all integers, and it has degree at most $N$. The claims now follow from \cref{lem:integer-valued-coefficients} applied to $p^{L_p}q_{ij}$ and to $b_i$.
\end{proof}

\begin{lemma}[Newton interpolation]\label{lem:newton}
Let $p$ be a prime and $\phi\ge0$ an integer with $v_p(\lambda/(1-\lambda))\ge-\phi$ and $v_p(1/(1-\lambda))\ge-\phi$. Then for every polynomial $g$ of degree at most $D$,
\[
v_p\bigl(G(g)\bigr)\ \ge\ \min_{0\le u\le D}v_p\bigl(g(u)\bigr)-(D+1)\phi .
\]
\end{lemma}

\begin{proof}
Write $g=\sum_{k\le D}\Delta^kg(0)\,b_k$ with $\Delta^kg(0)=\sum_{u\le k}(-1)^{k-u}\binom ku g(u)$. Then $G(b_0)=\lambda/(1-\lambda)$ and $G(b_k)=(\lambda/(1-\lambda))^k(1-\lambda)^{-1}$ for $k\ge1$, both of valuation at least $-(k+1)\phi$.
\end{proof}

For the four arguments one may take $\phi=0$ at every prime except $\pP$, where one may take $\phi=f$ (\cref{tab:parameters}). Indeed $\lambda/(1-\lambda)=2\epsilon/(b-2\epsilon)$ and $1/(1-\lambda)=b/(b-2\epsilon)$, and $b-2\epsilon=\pP^f$.

\subsection{Content}\label{sec:content}
Let $F\in\Q[X]$ be nonzero, choose an integer $D\ge1$ with $DF\in\Z[X]$, and let $g\ge1$ be the greatest common divisor of the coefficients of $DF$. The number $\cont(F)=g/D$ does not depend on $D$, and $F/\cont(F)$ is a primitive polynomial with integer coefficients. For every prime $p$ one has $v_p(\cont F)=\vg p(F)$, so that
\begin{equation}\label{eq:content-sum}
\log\cont(F)=\sum_p\vg p(F)\log p ,
\end{equation}
a finite sum. The Gauss valuation satisfies $\vg p(F+F')\ge\min(\vg p(F),\vg p(F'))$ and $\vg p(FF')=\vg p(F)+\vg p(F')$. In particular, if every entry of an $h\times h$ matrix with entries in $\Q[X]$ in row $i$ and column $j$ has Gauss valuation at least $\mu_i+\mu'_j$, then its determinant has Gauss valuation at least $\sum_i(\mu_i+\mu'_i)$. The arithmetic estimates below are lower bounds for $\vg p(\Qt_n)$; they hold for all $n$, trivially when $\Qt_n=0$.

\section{Arithmetic I: uniform bounds and the exceptional primes}\label{sec:exceptional}

This section bounds $\vg p(\Qt_n)$ from below at the three exceptional primes, and at all other primes by a crude uniform bound. The exceptional primes contribute a fixed multiple of $h^2$ to $\log\cont(\Qt_n)$ (\cref{cor:exceptional}); the crude bound is used for the small primes $p<\delta h$ and the large primes $p>K$ in the proof of \cref{thm:arith-general}. The statements are candidates in the sense of \cref{sec:intro-status}; their finite forms were checked on exact data (\cref{app:numerics}). The arguments are those of \cite[Sections~4.1 and 9.1--9.7]{Litsey2026Li2} for Tang's family, adapted to general parameters $(H,B,C)$: \cref{prop:P} corresponds to Lemma~4.2 there, and \cref{lem:reflection,prop:Q} to Lemma~9.9 and Proposition~9.8.

Let $\Ex=\{2,\pP,\pQ\}$ be the set of exceptional primes of \cref{tab:parameters}. A prime $p\notin\Ex$ is called ordinary; for such $p$ both $\lambda$ and $1-\lambda$ are $p$-adic units. All statements of this section hold for every $n\ge1$. We use the parameter conditions
\begin{equation}\label{eq:parameter-conditions}
h\le K,\qquad K-A\ge h,\qquad 0\le d\le 2,
\end{equation}
which hold for the four triples of \cref{tab:parameters}.

\begin{proposition}[ordinary primes, uniform bound]\label{prop:generic}
Let $p$ be an ordinary prime and $L_p$ as in \eqref{eq:N-L}. Then $\vg p(\Qt_n)\ge-2hL_p$. If moreover $p>K$, then $\vg p(\Qt_n)\ge0$.
\end{proposition}

\begin{proof}
We bound the entries $\Lambda_X(R_{ij})$ of the matrix of \cref{lem:normalization}, written as in \eqref{eq:functional-explicit}. By \cref{lem:integral-valued}(c) with $s=0,1$, the polynomial $(xq_{ij})'=q_{ij}+xq_{ij}'$ takes values of valuation at least $-2L_p$ at all integers; by \cref{lem:newton} with $\phi=0$, $v_p(G((xq_{ij})'))\ge-2L_p$. The residues $c^{(ij)}_l$ are integers, $\lambda^{-l}$ is a unit, and $v_p(\tau_l(\lambda))\ge-2\lfloor\log l/\log p\rfloor\ge-2L_p$ because $l\le K\le N$. Every entry therefore has Gauss valuation at least $-2L_p$, and the determinant at least $-2hL_p$.

Suppose $p>K$. Since $h\le K<p$ by \eqref{eq:parameter-conditions}, the polynomials $b_i$ ($i<h$) have coefficients in $\Zp$, and so has the numerator $K!\binom{x+A}A^3b_ib_j$ of $R_{ij}$. Division by the monic polynomial $D_K$ preserves this, so $q_{ij}\in\Zp[x]$, and $G((xq_{ij})')\in\Zp$ by \eqref{eq:moments}, the denominators there being powers of $1-\lambda$. Each $\tau_l$ with $l\le K<p$ lies in $\Zp$. All entries are $p$-integral, and $\vg p(\Qt_n)\ge0$.
\end{proof}

\begin{proposition}[the prime $2$]\label{prop:two}
$\vg2(\Qt_n)\ \ge\ -\sum_{k=0}^{h-1}(K-k)-2hL_2\ =\ -\bigl(\kappa-\tfrac12\bigr)h^2-\tfrac h2-2hL_2 .$
\end{proposition}

\begin{proof}
At $p=2$ we may take $\phi=0$ in \cref{lem:newton}, and each $\tau_l(\lambda)$ is $2$-integral because $v_2(\lambda^k/k^2)=k-2v_2(k)\ge0$. By \cref{lem:integral-valued}(a), the matrix of \cref{lem:normalization} is
\[
\Pi+\sum_{l=A+1}^{K}s_l\,\mathbf v_l\mathbf v_l^{\mathsf T},\qquad \Pi_{ij}=G\bigl((xq_{ij})'\bigr),\quad \mathbf v_l=\bigl(b_i(-l)\bigr)_{i<h}\in\Z^h,\quad s_l=l\gamma_l\lambda^{-l}\bigl(X-\tau_l(\lambda)\bigr),
\]
with $v_2(\Pi_{ij})\ge-2L_2$ as in \cref{prop:generic} and $\vg2(s_l)\ge-l$. Write $\Pi=\sum_{j<h}\Pi_{\bullet j}\mathbf e_j^{\mathsf T}$, where $\Pi_{\bullet j}$ is the $j$th column. The matrix is then a sum of rank-one terms of two kinds: $\Pi_{\bullet j}\mathbf e_j^{\mathsf T}$, of weight at least $-2L_2$, and $(s_l\mathbf v_l)\mathbf v_l^{\mathsf T}$, of weight at least $-l$. By the slot lemma (\cref{lem:slots}) the determinant has Gauss valuation at least the minimum, over $h$ distinct terms, of the sum of their weights. Since $K-A\ge h$ by \eqref{eq:parameter-conditions}, this minimum is at least $-\sum_{k<h}(K-k)-2hL_2$.
\end{proof}

\begin{proposition}[the prime $\pP$]\label{prop:P}
$\vg\pP(\Qt_n)\ \ge\ -2hL_{\pP}-f\sum_{i=0}^{h-1}\bigl(2i+3A-K+1\bigr)_+\ \ge\ -\tfrac f4(2-d)^2h^2-O(h\log h).$
\end{proposition}

\begin{proof}
At $\pP$ the argument $\lambda$ is a unit, so the pole terms are bounded as in \cref{prop:generic}. In \cref{lem:newton} we take $\phi=f$, and $\deg(xq_{ij})'\le\deg q_{ij}\le i+j+3A-K$ by \cref{lem:integral-valued}(b). Hence, with $\psi(m)=(m+3A-K+1)_+$, every entry satisfies $\vg\pP(\Lambda_X(R_{ij}))\ge-2L_{\pP}-f\psi(i+j)$ (if $q_{ij}=0$ the polynomial part contributes nothing). Expanding over permutations $\pi$,
\[
\vg\pP(\Qt_n)\ge-2hL_{\pP}-f\max_\pi\sum_{i<h}\psi\bigl(i+\pi(i)\bigr)\ge-2hL_{\pP}-f\sum_{i<h}\psi(2i),
\]
because $\psi$ is convex and $\psi(i+j)\le\frac12(\psi(2i)+\psi(2j))$. Finally $\sum_{i<h}(2i-dh+1)_+=\frac14(2-d)^2h^2+O(h)$ because $0\le d\le2$ by \eqref{eq:parameter-conditions}, and $L_{\pP}=O(\log h)$.
\end{proof}

The prime $\pQ$ divides the denominator of $\lambda$, so the factors $\lambda^{-l}$ in \cref{def:functional} are divisible by large powers of $\pQ$; but so are the denominators of $\tau_l(\lambda)$. The following identity separates the two effects.

\begin{lemma}[reflection formula]\label{lem:reflection}
Let $R\in\cR$ have polynomial part $q$ and poles in $\{-1,\dots,-K\}$, and for $1\le m\le K$ write its Laurent expansion at $-m$ as
\[
R(x)=\frac{c_m}{x+m}+r_{m,0}+r_{m,1}(x+m)+O\bigl((x+m)^2\bigr).
\]
Put $C_K=\sum_{k=1}^K\lambda^{-k}k^{-2}$ and $g=(xq)'$. Then
\begin{equation}\label{eq:reflection}
\Lambda_X(R)=-R(0)+\sum_{m=1}^K\lambda^{-m}\bigl(mc_m(X+C_K)-r_{m,0}+mr_{m,1}\bigr)+\cT(R),
\end{equation}
where
\[
\cT(R)=\lambda^{-K-1}G\bigl(g(x-K-1)\bigr)-\sum_{l=1}^Klc_l\sum_{k=K-l+1}^K\frac{\lambda^{-l-k}}{k^2}.
\]
\end{lemma}

\begin{proof}
For every polynomial $g$, shifting the summation index in the convergent series defining $G(g(x-K-1))$ gives
\begin{equation}\label{eq:shift}
G(g)=-\sum_{m=0}^K\lambda^{-m}g(-m)+\lambda^{-K-1}G\bigl(g(x-K-1)\bigr).
\end{equation}
Since $xR=xq+\sum_lc_l-\sum_l lc_l(x+l)^{-1}$, we have $g=(xR)'-\sum_llc_l(x+l)^{-2}$. Comparing Laurent expansions at $-m$ gives
\[
g(-m)=r_{m,0}-mr_{m,1}-\sum_{l\ne m}\frac{lc_l}{(l-m)^2}\quad(1\le m\le K),\qquad g(0)=R(0)-\sum_l\frac{c_l}l .
\]
Insert these values into \eqref{eq:shift} and add the pole terms $\sum_llc_l\lambda^{-l}(X-\tau_l(\lambda))$ of \eqref{eq:functional-explicit}. For fixed $l$, the terms $m<l$ of the double sum $\sum_m\lambda^{-m}\sum_{l\ne m}lc_l(l-m)^{-2}$ equal $lc_l\sum_{k=1}^{l-1}\lambda^{k-l}k^{-2}$; together with $-lc_l\lambda^{-l}\tau_l(\lambda)=-lc_l\sum_{k=1}^{l}\lambda^{k-l}k^{-2}$ they leave $-c_l/l$, which cancels the corresponding term of $-g(0)$. The terms $m>l$ equal $lc_l\lambda^{-l}\sum_{k=1}^{K-l}\lambda^{-k}k^{-2}=lc_l\lambda^{-l}\bigl(C_K-\sum_{k=K-l+1}^K\lambda^{-k}k^{-2}\bigr)$. Collecting terms gives \eqref{eq:reflection}.
\end{proof}

\begin{proposition}[the prime $\pQ$]\label{prop:Q}
$\vg\pQ(\Qt_n)\ \ge\ e\sum_{k=0}^{h-2}\Bigl(A+1+\Bigl\lfloor\frac k3\Bigr\rfloor\Bigr)-2hL_{\pQ}\ =\ e\bigl(a+\tfrac16\bigr)h^2-O(h\log h).$
\end{proposition}

\begin{proof}
Write $v=v_{\pQ}$ and $L=L_{\pQ}$. We have $v(\lambda^{-1})=e$, $v(\lambda/(1-\lambda))=0$ and $v(1/(1-\lambda))=e$, so $\phi=0$ in \cref{lem:newton}. The constant $C_K$ is $\pQ$-integral, because $ek-2v(k)\ge0$ for $k\ge1$.

Apply \cref{lem:reflection} to $R=R_{ij}$. First, $R_{ij}(0)=S_nW_n(0)b_i(0)b_j(0)=\ind_{i=j=0}$. Second, for $m\le A$ the function $R_{ij}$ vanishes to order two at $-m$, so $c_m=r_{m,0}=r_{m,1}=0$. Third, for $A<m\le K$ write $S_nW_n(x)=\chi_m(y)/y$ with $y=x+m$; here
\[
\chi_m(y)=K!\prod_{a=1}^A\Bigl(\frac{a-m+y}{a}\Bigr)^3\prod_{\substack{1\le l\le K\\ l\ne m}}(l-m+y)^{-1},
\]
so $\chi_m(0)=\gamma_m\in\Z$, and since $|a-m|,|l-m|\le K\le N$, every coefficient satisfies $v([y^u]\chi_m)\ge-uL$. Let $\beta^{(m)}_{i,s}=[y^s]b_i(y-m)$, of valuation at least $-sL$ by \cref{lem:integral-valued}(c). Then $c_m$, $r_{m,0}$, $r_{m,1}$ are the coefficients of $y^0,y^1,y^2$ in $\chi_m(y)b_i(y-m)b_j(y-m)$, and
\[
mc_m(X+C_K)-r_{m,0}+mr_{m,1}=\sum_{s+s'\le2}\beta^{(m)}_{i,s}\beta^{(m)}_{j,s'}\,\Omega^{(m)}_{s+s'},
\]
with $\Omega^{(m)}_0=m\chi_{m,0}(X+C_K)-\chi_{m,1}+m\chi_{m,2}$, $\Omega^{(m)}_1=-\chi_{m,0}+m\chi_{m,1}$, $\Omega^{(m)}_2=m\chi_{m,0}$, and $\chi_{m,u}=[y^u]\chi_m$. Thus $\vg{}(\Omega^{(m)}_{u})\ge-(2-u)L$. Finally, the tail $\cT(R_{ij})$ has valuation at least $e(K+1)-2L$: its first term because $v(G(g(x-K-1)))\ge-2L$ as in \cref{prop:generic}, its second because $v(\lambda^{-l-k}k^{-2})\ge e(l+k)-2L$ with $l+k\ge K+1$.

The matrix $[\Lambda_X(R_{ij})]$ is therefore the sum of the following rank-one terms:
\begin{itemize}
\item one constant term $-\mathbf e_0\mathbf e_0^{\mathsf T}$, of weight $0$;
\item for $A<m\le K$ and $s\in\{0,1,2\}$, the term $\mathbf u_{m,s}\mathbf w_{m,s}^{\mathsf T}$ with $\mathbf u_{m,s}=(\beta^{(m)}_{i,s})_i$ and $\mathbf w_{m,s}=(\lambda^{-m}\sum_{s'\le2-s}\Omega^{(m)}_{s+s'}\beta^{(m)}_{j,s'})_j$, of weight at least $-sL+em-(2-s)L=em-2L$;
\item for $i<h$, the term $\mathbf e_i\,\cT_{i\bullet}$ formed with the $i$th row of the tail matrix, of weight at least $e(K+1)-2L$.
\end{itemize}
By the slot lemma (\cref{lem:slots}), $\vg\pQ(\Qt_n)$ is at least the minimal total weight of $h$ distinct terms. Ordering the non-constant terms by weight, the $k$th of them ($k=0,1,\dots,h-1$) has weight at least $e(A+1+\lfloor k/3\rfloor)-2L$, since there are three terms for each $m$ and $K-A\ge h$. A choice of $h$ terms that contains the constant term has total weight at least $e\sum_{k\le h-2}(A+1+\lfloor k/3\rfloor)-2(h-1)L$. A choice of $h$ non-constant terms has total weight at least $e\sum_{k\le h-1}(A+1+\lfloor k/3\rfloor)-2hL$, which exceeds the stated bound by $e(A+1+\lfloor(h-1)/3\rfloor)\ge0$. In both cases the total is at least the stated bound; the constant term cannot simply be assumed to be selected, because for $A=0$ the smallest non-constant weights $e-2L$ are negative. The asymptotic form follows from $\sum_{k\le h-2}\lfloor k/3\rfloor=h^2/6+O(h)$.
\end{proof}

\begin{corollary}\label{cor:exceptional}
As $n\to\infty$,
\[
\sum_{p\in\Ex}\vg p(\Qt_n)\log p\ \ge\ \eps_{\mathrm{exc}}h^2-O(h\log h),
\]
where
\[
\eps_{\mathrm{exc}}=-\bigl(\kappa-\tfrac12\bigr)\log2-\tfrac f4(2-d)^2\log\pP+e\bigl(a+\tfrac16\bigr)\log\pQ .
\]
Numerically $\eps_{\mathrm{exc}}=-0.552315$, $-0.439472$, $-0.550398$ and $-0.598268$ for $\lambda=2/11$, $2/9$, $-2/9$ and $-2/7$ respectively (rounded to six decimals).
\end{corollary}

\begin{proof}
Combine \cref{prop:two,prop:P,prop:Q}.
\end{proof}

\section{Arithmetic II: ordinary primes and residue classes}\label{sec:ordinary}

This section proves the two local lower bounds for $\vg p(Q_n)$ at an ordinary prime $p>\sqrt K$, the jet bound (\cref{prop:jet}) and the pole bound (\cref{prop:pole}). They are the main source of the arithmetic gain, and \cref{sec:profile} integrates them over the primes $\delta h\le p\le K$; the ordinary primes $p\le\sqrt K$ are covered only by the uniform bound of \cref{prop:generic}. Note that the bounds of this section are stated for the determinant $Q_n$ itself, whereas those of \cref{sec:exceptional} are stated for $\Qt_n=(S_n^h/F_n)Q_n$; the valuation of the factor $S_n^h/F_n$ is added in \cref{prop:profile-uniform}. The same distribution over residue classes is used again in \cref{sec:nonvanishing}. The statements are candidates in the sense of \cref{sec:intro-status}.

The method is that of \cite[Sections~9.8 and~9.9]{Litsey2026Li2}, where it is attributed to Tang: the residue classes modulo $p$, the counts $a_\sigma$, $c_\sigma$ and the exponent $\nu_\sigma$ of \cref{lem:local-integrality} appear there for Tang's family, with a bound for the entries (Proposition~9.12 there) and a determinant lemma for point masses at congruent nodes (Lemma~9.5 there); the bound for the primes $n<p\le4n$ (Proposition~9.15 there) already contains the shift by one unit for the poles divisible by $p$. \Cref{prop:jet,prop:pole} state the outcome for general parameters as sums over slot multisets. Our proof of the local integrality, through finite geometric averages and an affine change of the indeterminate, is different from the proof in \cite{Litsey2026Li2}; we have not compared the two in detail.

Throughout this section $p$ is an ordinary prime with $K<p^2$, and $n$ is arbitrary. Then $p\ge3$, and $\lambda$, $1-\lambda$ are $p$-adic units. Put $\lambda_p=\lambda^p$. By Fermat's theorem $\lambda_p\equiv\lambda\pmod p$, so $\lambda_p$ and $1-\lambda_p$ are also units; more generally $\lambda_p^M\equiv\lambda\pmod p$ for every power $M$ of $p$. For $0\le\sigma<p$ let
\[
a_\sigma=\#\{1\le l\le A:\ l\equiv\sigma\ (\mathrm{mod}\ p)\},\qquad c_\sigma=\#\{1\le l\le K:\ l\equiv\sigma\ (\mathrm{mod}\ p)\},
\]
and $w_\sigma=3a_\sigma-c_\sigma-1$.
The idea is to split the sum $\sum_{x\ge1}\lambda^x f(x)$ according to the class of $-x$ modulo $p$, writing $x=pu-\sigma$ with $u\ge1$ and $0\le\sigma<p$. On each class the weight $W_n$ is, up to a $p$-adic unit, the power $p^{3a_\sigma-c_\sigma}$ times a rational function in $u$ with small poles; this is the source of the valuations $w_\sigma$. Since $\lambda^x$ is not $p$-adically continuous in $x$, we work with finite normalized geometric sums, and we transport the indeterminate $X$ by an affine change of variable.

\subsection{Finite geometric averages}\label{sec:finite-averages}
For a rational number $v$ with $0<|v|<1$, an integer $M\ge1$ and a function $f$ defined on the positive integers put
\[
T_{M,v}(f)=\frac1{1-v^M}\sum_{u=1}^Mv^uf(u).
\]

\begin{lemma}\label{lem:finite-moments}
Let $v$ be a rational number with $0<|v|<1$ such that $v$ and $1-v$ are $p$-adic units, let $\mu\ge1$ and $M=p^\mu$. Then $1-v^M$ is a unit, and:
\begin{enumerate}[(a)]
\item if $f$ takes values in $\Zp$ at the positive integers, then $T_{M,v}(f)\in\Zp$;
\item if $g\in\Zp[u]$, then $G_v(g)\in\Zp$ and $T_{M,v}(g)\equiv G_v(g)\pmod{p^\mu}$.
\end{enumerate}
\end{lemma}

\begin{proof}
We have $v^M\equiv v\pmod p$, so $1-v^M$ is a unit, and (a) follows. For (b), $G_v(g)=\sum_k\Delta^kg(0)\,G_v(b_k)$ with $\Delta^kg(0)\in\Zp$ and $G_v(b_k)\in\Zp$ by the formulas in the proof of \cref{lem:newton}. Splitting the convergent series $G_v(g)$ at $u=M$ gives $\sum_{u=1}^Mv^ug(u)=G_v(g)-v^MG_v(g(u+M))$, whence
\[
T_{M,v}(g)-G_v(g)=\frac{v^M}{1-v^M}\bigl(G_v(g)-G_v(g(u+M))\bigr).
\]
By Taylor's formula $g(u+M)-g(u)=\sum_{k\ge1}M^k\,g^{(k)}(u)/k!$, and $g^{(k)}/k!\in\Zp[u]$; so $g(u+M)-g(u)\in p^\mu\Zp[u]$ and the difference is divisible by $p^\mu$.
\end{proof}

The functional $\Lambda_X$ acts on derivatives $f=(xR)'$, which are sums of a polynomial and of double poles $(x+l)^{-2}$. We regularize the double poles with an affine parameter. A \emph{function with designated double poles} is a function $f$ on the positive integers together with a decomposition $f=\varphi+\sum_{k\in S}d_k(u+k)^{-2}$, where $S\subset\Z_{\ge0}$ is finite, $d_k\in\Q$ and $\varphi$ is a function on the positive integers. For such $f$, a rational $v$ with $0<|v|<1$ and an indeterminate $Y$ put
\[
D_v(f)=\sum_{k\in S}d_kv^{-k},\qquad T_{M,v,Y}(f)=T_{M,v}(f)+D_v(f)\bigl(Y-T_{M,v}(u^{-2})\bigr)\in\Q[Y].
\]
This is linear in $f$ (with the designations added). Its point is that the dependence on $Y$ is concentrated in the designated poles, which are treated exactly as $\Lambda_X$ treats the poles of $(xR)'$.

\begin{lemma}[regularized double poles]\label{lem:pole-regularization}
Let $v$ be a rational number with $0<|v|<1$ and $M\ge1$. For $k\ge0$, with the single designated pole $(u+k)^{-2}$,
\[
T_{M,v,Y}\bigl((u+k)^{-2}\bigr)=v^{-k}\bigl(Y-\tau_k(v)\bigr)+\eta_{M,k}(v),
\]
where
\[
\eta_{M,k}(v)=\frac{v^{M-k}}{1-v^M}\sum_{j=1}^kv^j\Bigl(\frac1{(M+j)^2}-\frac1{j^2}\Bigr).
\]
If moreover $v$ and $1-v$ are $p$-adic units, $M=p^\mu$, and $\mu>v_p(j)$ for all $j\le k$, then $v_p(\eta_{M,k}(v))\ge \mu-3\max_{j\le k}v_p(j)$. In particular $v_p(\eta_{M,k}(v))\ge \mu$ if $k<p$.
\end{lemma}

\begin{proof}
Reindexing, $\sum_{u=1}^Mv^u(u+k)^{-2}=v^{-k}\bigl(\sum_{m=1}^Mv^mm^{-2}-\tau_k(v)+\sum_{j=1}^kv^{M+j}(M+j)^{-2}\bigr)$. Divide by $1-v^M$ and add $v^{-k}(Y-T_{M,v}(u^{-2}))$; the terms $T_{M,v}(u^{-2})$ cancel, and $-v^{-k}\tau_k(1-v^M)^{-1}=-v^{-k}\tau_k-v^{M-k}\tau_k(1-v^M)^{-1}$ produces the formula. For the valuation, $(M+j)^{-2}-j^{-2}=-M(M+2j)(M+j)^{-2}j^{-2}$, and $v_p(M+j)=v_p(M+2j)=v_p(j)$ when $\mu>v_p(j)$, since $p$ is odd.
\end{proof}

\subsection{Distribution over residue classes}\label{sec:distribution}
Let $R\in\cR$ have its poles in $\{-1,\dots,-K\}$ and put $f=(xR)'$, with designated double poles at all $-l$ ($1\le l\le K$). For $0\le\sigma<p$ let $f_\sigma(u)=f(pu-\sigma)$ ($u\ge1$). A pole $(x+l)^{-2}$ of $f$ with $l\equiv\sigma\pmod p$ becomes $p^{-2}(u+k)^{-2}$ with $k=(l-\sigma)/p$; since $l\le K<p^2$ we have $0\le k<p$. We designate exactly these poles in $f_\sigma$; the others, $(pu+l-\sigma)^{-2}$ with $p\nmid l-\sigma$, are part of the regular term and take values in $\Zp$ at integers. For $\mu\ge1$ put $M=p^\mu$ and
\[
L^*_\mu=\sum_{\sigma=1}^{p-1}\lambda^{-\sigma}\,T_{M,\lambda_p}\bigl((pu-\sigma)^{-2}\bigr),\qquad Y_\mu=p^2\bigl(X-L^*_\mu\bigr).
\]
By \cref{lem:finite-moments}(a), $L^*_\mu\in\Zp$, so both coefficients of the affine polynomial $Y_\mu$ lie in $p^2\Zp$; in particular $\vg p(Y_\mu)\ge2$, and $Y_\mu\equiv0\pmod{p^2}$ after specializing $X$ to any $r\in\Zp$.

\begin{proposition}[distribution]\label{prop:distribution}
With this notation there is a constant $c_R$, independent of $\mu$, such that
\[
\vg p\Bigl(\Lambda_X(R)-\sum_{\sigma=0}^{p-1}\lambda^{-\sigma}\,T_{M,\lambda_p,Y_\mu}(f_\sigma)\Bigr)\ \ge\ \mu-c_R .
\]
\end{proposition}

\begin{proof}
We first prove the exact identity
\begin{equation}\label{eq:distribution-exact}
T_{pM,\lambda,X}(f)=\sum_{\sigma=0}^{p-1}\lambda^{-\sigma}\,T_{M,\lambda_p,Y_\mu}(f_\sigma).
\end{equation}
The map $(u,\sigma)\mapsto pu-\sigma$ is a bijection from $\{1,\dots,M\}\times\{0,\dots,p-1\}$ onto $\{1,\dots,pM\}$, and $\lambda^{pu-\sigma}=\lambda^{-\sigma}\lambda_p^u$, $1-\lambda^{pM}=1-\lambda_p^M$. Hence $T_{pM,\lambda}(\varphi)=\sum_\sigma\lambda^{-\sigma}T_{M,\lambda_p}(\varphi(pu-\sigma))$ for every function $\varphi$; applied to $f$ and to $x^{-2}$ this gives
\[
T_{pM,\lambda}(f)=\sum_\sigma\lambda^{-\sigma}T_{M,\lambda_p}(f_\sigma),\qquad T_{pM,\lambda}(x^{-2})=p^{-2}T_{M,\lambda_p}(u^{-2})+L^*_\mu .
\]
If $f$ has the designated pole $d_l(x+l)^{-2}$, then $f_\sigma$ has the designated pole $p^{-2}d_l(u+k)^{-2}$ for $\sigma\equiv l$, so $D_{\lambda_p}(f_\sigma)=p^{-2}\sum_{l\equiv\sigma}d_l\lambda^{-(l-\sigma)}$ and $\sum_\sigma\lambda^{-\sigma}D_{\lambda_p}(f_\sigma)=p^{-2}D_\lambda(f)$. Therefore the right side of \eqref{eq:distribution-exact} equals
\[
T_{pM,\lambda}(f)+p^{-2}D_\lambda(f)\bigl(Y_\mu-T_{M,\lambda_p}(u^{-2})\bigr)=T_{pM,\lambda}(f)+D_\lambda(f)\bigl(X-L^*_\mu-p^{-2}T_{M,\lambda_p}(u^{-2})\bigr),
\]
which is $T_{pM,\lambda,X}(f)$.

It remains to compare $T_{pM,\lambda,X}(f)$ with $\Lambda_X(R)$. Write $R=q_R+\sum_lc_l(x+l)^{-1}$, so that $f=(xq_R)'+\sum_l lc_l(x+l)^{-2}$. By linearity and \cref{lem:pole-regularization} with $v=\lambda$,
\[
T_{pM,\lambda,X}(f)=T_{pM,\lambda}\bigl((xq_R)'\bigr)+\sum_llc_l\bigl(\lambda^{-l}(X-\tau_l(\lambda))+\eta_{pM,l}(\lambda)\bigr).
\]
Comparing with \eqref{eq:functional-explicit}, the difference from $\Lambda_X(R)$ is $T_{pM,\lambda}(g)-G(g)$ with $g=(xq_R)'$, plus $\sum_llc_l\eta_{pM,l}(\lambda)$. Choose $c\ge0$ with $p^cg\in\Zp[x]$ and $p^cc_l\in\Zp$; by \cref{lem:finite-moments}(b) and \cref{lem:pole-regularization} (where $v_p(j)\le1$ for $j\le K<p^2$) the difference has valuation at least $\mu+1-c-3$.
\end{proof}

We apply this to $R=W_nE\,E'$ with polynomials $E,E'$. Since $f(pu-\sigma)=p^{-1}\frac{\mathrm d}{\mathrm du}\bigl[(pu-\sigma)R(pu-\sigma)\bigr]$, it is natural to define, for $0\le\sigma<p$ and $\Phi\in\Zp[u]$,
\begin{equation}\label{eq:class-term}
\Theta_\sigma(\Phi)=\frac{\lambda^{-\sigma}}{p}\,T_{M,\lambda_p,Y_\mu}\Bigl(\frac{\mathrm d}{\mathrm du}\bigl[(pu-\sigma)W_n(pu-\sigma)\,\Phi(u)\bigr]\Bigr),
\end{equation}
with the designated poles described above (the dependence on $\mu$ is suppressed). \Cref{prop:distribution} says that $\Lambda_X(W_nEE')\equiv\sum_\sigma\Theta_\sigma\bigl(E(pu-\sigma)E'(pu-\sigma)\bigr)$ modulo $p^{\mu-c}$, coefficientwise in $X$, with $c$ independent of $\mu$.

\subsection{Local integrality}\label{sec:local-integrality}
For $\Phi\in\Zp[u]$ let $o(\Phi)$ be the minimum valuation of its coefficients.

\begin{lemma}[local integrality]\label{lem:local-integrality}
For $0\le\sigma<p$, $\mu\ge1$ and $\Phi\in\Zp[u]$,
\[
\vg p\bigl(\Theta_\sigma(\Phi)\bigr)\ \ge\ \nu_\sigma+o(\Phi),\qquad \nu_\sigma=3a_\sigma-c_\sigma-1+\ind_{\sigma=0}.
\]
\end{lemma}

\begin{proof}
For $1\le a\le K$, the factor $pu-\sigma+a$ equals $p(u+k)$ with $k=(a-\sigma)/p\in\{0,\dots,p-1\}$ when $a\equiv\sigma$, and is a polynomial in $u$ with unit constant term and all other coefficients divisible by $p$ otherwise. Collecting the factors of $W_n$, and the extra factor $pu$ when $\sigma=0$,
\[
(pu-\sigma)W_n(pu-\sigma)\Phi(u)=p^{3a_\sigma-c_\sigma+\ind_{\sigma=0}}\ \frac{\Phi_1(u)}{B(u)\,D(u)},
\]
where $B(u)=\prod_{k\in S}(u+k)$ is a product of distinct monic factors with $S\subset\{0,\dots,p-1\}$ (the local poles of $W_n$ that remain after cancellation with the local zeros; since $A\le K$ by \eqref{eq:parameter-conditions}, every local zero $l\le A$ is also a local pole, so the cancellation leaves a polynomial numerator), $D\in\Zp[u]$ is the product of the non-local denominator factors, so $D\equiv D(0)\not\equiv0\pmod p$ coefficientwise, and $\Phi_1\in\Zp[u]$ with $o(\Phi_1)\ge o(\Phi)$. In particular $D(u)\in\Zp^\times$ for every integer $u$.

For $k\in S$ the residue of $\Phi_1/(BD)$ at $-k$ is $\varrho_k=\Phi_1(-k)/\bigl(D(-k)\prod_{k'\in S\setminus\{k\}}(k'-k)\bigr)$, of valuation at least $o(\Phi)$, because $0<|k'-k|<p$. The polynomial $\Phi_1-D\sum_k\varrho_k\prod_{k'\ne k}(u+k')$ vanishes at every $-k$, $k\in S$, hence is divisible by the monic polynomial $B$ in $\Zp[u]$:
\[
\frac{\Phi_1}{BD}=\sum_{k\in S}\frac{\varrho_k}{u+k}+\frac{\Phi_0}{D},\qquad \Phi_0\in\Zp[u],\ o(\Phi_0)\ge o(\Phi).
\]
Differentiating, the function in \eqref{eq:class-term} is $p^{3a_\sigma-c_\sigma+\ind_{\sigma=0}}$ times $-\sum_k\varrho_k(u+k)^{-2}+(\Phi_0'D-\Phi_0D')D^{-2}$, and the designated poles are exactly the $(u+k)^{-2}$, $k\in S$. The regular part takes values of valuation at least $o(\Phi)$ at the integers, so its average $T_{M,\lambda_p}$ does too by \cref{lem:finite-moments}(a). Each designated pole contributes $-\varrho_k\bigl(\lambda_p^{-k}(Y_\mu-\tau_k(\lambda_p))+\eta_{M,k}(\lambda_p)\bigr)$ by \cref{lem:pole-regularization}; here $k<p$, so $\tau_k(\lambda_p)\in\Zp$ and $v_p(\eta_{M,k}(\lambda_p))\ge \mu\ge1$, and $\vg p(Y_\mu)\ge0$. Dividing by $p$ gives the claim.
\end{proof}

\subsection{The jet bound}\label{sec:jet-bound}
For a finite or infinite multiset $\cS$ of integers, bounded below, and $m\ge0$, let $\Sigma_m(\cS)$ be the sum of the $m$ smallest elements of $\cS$, and $\Sigma^-_m(\cS)$ the sum of the negative elements among them.

\begin{proposition}[jet bound]\label{prop:jet}
Let $\cS^{\mathrm{jet}}_p$ be the multiset $\{w_\sigma+2j:\ 0\le\sigma<p,\ j\ge0\}$. Then $\vg p(Q_n)\ge\cJ_p:=\Sigma_h(\cS^{\mathrm{jet}}_p)$.
\end{proposition}

\begin{proof}
Write $\mathsf M=[\Lambda_X(W_nx^{i+j})]_{i,j<h}$, so that $Q_n=\det\mathsf M$. Expanding $(pu-\sigma)^i=\sum_{s\le i}\binom is p^s(-\sigma)^{i-s}u^s$ and using \cref{prop:distribution} and the linearity of $\Theta_\sigma$, we get $\mathsf M\equiv\mathcal A\,\mathcal H\,\mathcal A^{\mathsf T}$ modulo $p^{\mu-c}$, coefficientwise, where
\begin{itemize}
\item $\mathcal A$ is the $h\times ph$ matrix with entries $\mathcal A_{i,(\sigma,s)}=\binom is p^s(-\sigma)^{i-s}$ ($0\le\sigma<p$, $0\le s<h$), whose column $(\sigma,s)$ has valuation at least $s$;
\item $\mathcal H$ is block diagonal with blocks $\mathcal H^{(\sigma)}=[\Theta_\sigma(u^{s+s'})]_{s,s'<h}$, whose entries have Gauss valuation at least $\nu_\sigma\ge w_\sigma$ by \cref{lem:local-integrality}.
\end{itemize}
By the Cauchy--Binet formula applied twice,
\[
\det(\mathcal A\mathcal H\mathcal A^{\mathsf T})=\sum_{I,J}\det\mathcal A_{\bullet I}\,\det\mathcal H_{IJ}\,\det\mathcal A_{\bullet J},
\]
the sum running over pairs of $h$-element sets $I,J$ of column indices. Since $\mathcal H$ is block diagonal, $\det\mathcal H_{IJ}=0$ unless $I$ and $J$ meet every block $\sigma$ in the same number $m_\sigma$ of indices, and then $\vg p(\det\mathcal H_{IJ})\ge\sum_\sigma m_\sigma w_\sigma$. Hence the term $(I,J)$ has Gauss valuation at least
\[
\sum_{(\sigma,s)\in I}\bigl(s+\tfrac12w_\sigma\bigr)+\sum_{(\sigma,s)\in J}\bigl(s+\tfrac12w_\sigma\bigr)\ \ge\ 2\min_{|I|=h}\sum_{(\sigma,s)\in I}\bigl(s+\tfrac12w_\sigma\bigr)\ \ge\ \Sigma_h(\cS^{\mathrm{jet}}_p).
\]
All entries of $\mathcal A\mathcal H\mathcal A^{\mathsf T}$ have Gauss valuation at least $\min_\sigma w_\sigma\ge-K-1$, independently of $\mu$. Therefore $\vg p(\det\mathsf M-\det\mathcal A\mathcal H\mathcal A^{\mathsf T})\ge \mu-c-(h-1)(K+1)$, and letting $\mu\to\infty$ proves the claim.
\end{proof}

\subsection{The pole bound}\label{sec:pole-bound}
The jet bound ignores the fact that the local poles of a class are evaluated at distinct points $l$. The second bound uses this through Vandermonde factors.

\begin{proposition}[pole bound]\label{prop:pole}
Assume moreover $p\le K$. Let $\cS^{\mathrm{pole}}_p$ be the multiset consisting of
\begin{align*}
&w_\sigma+2j+2\cdot\ind_{a_\sigma=0,\ j=c_\sigma-a_\sigma-1}&&(1\le\sigma<p,\ 0\le j<c_\sigma-a_\sigma),\\
&w_0+1+2j&&(0\le j<c_0-a_0).
\end{align*}
Then $\vg p(Q_n)\ge\cP_p:=\Sigma^-_h(\cS^{\mathrm{pole}}_p)$.
\end{proposition}

\begin{proof}
The polynomial part $\Pi_s$ of $W_nx^s$ has integer coefficients, since $D_K$ is monic, so $\Lambda_X(\Pi_s)=G((x\Pi_s)')\in\Zp$, where $G=G_\lambda$ is the geometric sum of \cref{sec:moments}. Let $\mathsf M=[\Lambda_X(W_nx^{i+j})]_{i,j<h}$ as in the proof of \cref{prop:jet}. The residue of $W_nx^s$ at $-l$, $A<l\le K$, is $\rho_l(-l)^s$ with $\rho_l=D_A(-l)^3/\prod_{m\ne l}(m-l)$. As $0<|a-l|,|m-l|<p^2$, each factor has valuation $1$ or $0$ according as it is divisible by $p$ or not, so $v_p(\rho_l)=3a_\sigma-(c_\sigma-1)=w_\sigma+2$ for $l\equiv\sigma$. Thus
\[
\mathsf M=\widehat\Pi+\sum_{l=A+1}^Ks_l\,\mathbf v_l\mathbf v_l^{\mathsf T},\qquad \widehat\Pi\in\Zp^{h\times h},\quad \mathbf v_l=\bigl((-l)^i\bigr)_{i<h},\quad s_l=l\rho_l\lambda^{-l}\bigl(X-\tau_l(\lambda)\bigr).
\]
We have $v_p(\tau_l(\lambda))\ge0$ if $l<p$ and $\ge-2$ if $p\le l<p^2$. Hence $\vg p(s_l)\ge w_\sigma+1$ if $\sigma=0$ (then $p\mid l$ and $l\ge p$), $\vg p(s_l)\ge w_\sigma+2$ if $l=\sigma\ne0$, and $\vg p(s_l)\ge w_\sigma$ otherwise. For $\sigma\ne0$, the pole $l=\sigma$ occurs exactly when $a_\sigma=0$: if $l=\sigma$ is a pole then $\sigma>A$, so no integer in $[1,A]$ is congruent to $\sigma$ and $a_\sigma=0$; conversely, if $a_\sigma=0$ then $\sigma>A$, and $\sigma<p\le K$ by the hypothesis $p\le K$, so $l=\sigma$ is a pole of the class. This is the only place where $p\le K$ is used.

Write $\mathsf M=[\widehat\Pi\,|\,V]\cdot\bigl[\begin{smallmatrix}I_h\\ SV^{\mathsf T}\end{smallmatrix}\bigr]$ with $V=[\mathbf v_l]_l$ and $S=\mathrm{diag}(s_l)$, and expand by Cauchy--Binet. Up to sign, the term indexed by a set $\mathsf L$ of poles and a set $I\subset\{0,\dots,h-1\}$ with $|I|+|\mathsf L|=h$ is
\[
\Bigl(\prod_{l\in\mathsf L}s_l\Bigr)\,\det\bigl[\widehat\Pi_{\bullet I}\,|\,V_{\bullet\mathsf L}\bigr]\,\det V_{I^c,\mathsf L}.
\]
Both determinants are values at the points $l\in\mathsf L$ of polynomials with $\Zp$-coefficients in variables $z_l$ that are alternating in these variables; by \cref{lem:alternating} each is divisible by $\prod_{l<l'}(z_l-z_{l'})$. Two poles in the same class differ by a multiple of $p$, so if $\mathsf L$ contains $k_\sigma$ poles of class $\sigma$, the two determinants together have valuation at least $\sum_\sigma k_\sigma(k_\sigma-1)$. For each class, $\sum_{l\in\mathsf L_\sigma}\vg p(s_l)+k_\sigma(k_\sigma-1)$ is at least the sum of the first $k_\sigma$ elements of the class-$\sigma$ part of $\cS^{\mathrm{pole}}_p$: for $\sigma=0$ because $k(k-1)=\sum_{j<k}2j$; for $\sigma\ne0$ for the same reason and because, when $a_\sigma=0$ and all $c_\sigma-a_\sigma$ poles of the class are selected, the pole $l=\sigma$ is among them and contributes the extra $2$. Within each class these elements are nondecreasing in $j$. The total is therefore at least the sum of a selection of at most $h$ elements of $\cS^{\mathrm{pole}}_p$, which is at least $\Sigma^-_h(\cS^{\mathrm{pole}}_p)$.
\end{proof}

\begin{remark}\label{rem:zero-class}
In the source of this argument the zero class was given the same slots as the nonzero classes. That version can fail by one unit: when $a_0=0$ and $c_0=1$, the only pole $l=p$ of the zero class has $\vg p(s_l)\ge w_0+1$, not $w_0+2$, and the exact computations of \cref{app:numerics} exhibit primes where the actual valuation of $Q_n$ lies exactly one unit below that version. The zero-class slots $w_0+1+2j$ of \cref{prop:pole} are correct, and the bound in this form holds at every prime and every $n$ tested. Since $\sum_{j<k}(w_0+1+2j)\ge\sum_{j<k}(w_0+2j)+2\cdot\ind_{k=c_0-a_0,\ a_0=0}-2$ for every $k$, \cref{prop:pole} implies the source version with a correction $-2$, which is the form in which the source bound is valid. Since only one class is affected, the change is $O(1)$ per prime and does not affect any asymptotic rate.
\end{remark}

\section{Arithmetic III: the limiting profile and the prime number theorem}\label{sec:profile}

This section turns the local bounds of \cref{sec:ordinary} into the arithmetic estimate \cref{thm:arith}. For primes $p$ of size comparable with $h$, the class counts $a_\sigma$, $c_\sigma$ and the valuation of the normalization $S_n^h/F_n$ depend, up to bounded errors, only on $x=p/h$ (\cref{prop:profile-uniform}); the prime number theorem then integrates the resulting profile (\cref{thm:arith-general}), and a finite computation evaluates the integral (\cref{prop:profile-certificate}). The first two steps are candidates and the third is certified, in the sense of \cref{sec:intro-status}. A profile of this kind, integrated with the prime number theorem, is used for Tang's family in \cite[Proposition~4.3 and Section~9.10]{Litsey2026Li2}.

\subsection{The profile}\label{sec:profile-definition}
For $0<x\le\kappa$ and $s\in\{a,\kappa,1\}$ put $q_s=\lfloor s/x\rfloor$ and $r_s=s-q_sx\in[0,x)$; when $a=0$ we have $q_a=r_a=0$. For $0<u<x$ define
\[
\alpha(u)=q_a+\ind_{u<r_a},\qquad \gamma(u)=q_\kappa+\ind_{u<r_\kappa},\qquad \omega(u)=3\alpha(u)-\gamma(u)-1 .
\]
The interval $(0,x)$ is the union of at most three intervals, cut at $r_a$ and $r_\kappa$, on which the pair $(\alpha,\gamma)$ is constant. Call the pairs that occur the \emph{types} at $x$, and let $\ell_\tau(x)>0$ be the total length of the set where the type is $\tau=(\alpha_\tau,\gamma_\tau)$; write $\omega_\tau=3\alpha_\tau-\gamma_\tau-1$.

\begin{definition}\label{def:profile}
The \emph{jet slots} at $x$ are the pairs $(\tau,j)$, $j\ge0$, with weight $\omega_\tau+2j$ and capacity $\ell_\tau(x)$. The \emph{pole slots} are the pairs $(\tau,j)$ with $0\le j<\gamma_\tau-\alpha_\tau$, weight $\omega_\tau+2j+2\cdot\ind_{\alpha_\tau=0,\ j=\gamma_\tau-\alpha_\tau-1}$ and capacity $\ell_\tau(x)$. Let
\begin{align*}
J(x)&=\min\Bigl\{\textstyle\sum\mathrm{weight}\cdot y\ :\ 0\le y\le\mathrm{capacity},\ \sum y=1\Bigr\}\quad\text{(jet slots)},\\
R(x)&=\min\Bigl\{\textstyle\sum\mathrm{weight}\cdot y\ :\ 0\le y\le\mathrm{capacity},\ \sum y\le1\Bigr\}\quad\text{(pole slots)},\\
\mathsf N(x)&=q_\kappa-3q_a-2\Bigl(\frac{x\,q_1(q_1-1)}2+q_1r_1\Bigr),\qquad \mathsf B(x)=\mathsf N(x)+\max\bigl(J(x),R(x)\bigr).
\end{align*}
\end{definition}

Only jet slots with $j\le q_1+2$ can carry mass in an optimal solution: each level has total capacity $\sum_\tau\ell_\tau(x)=x$, so the levels $j\le q_1$ have total capacity $(q_1+1)x>1$; the weights $\omega_\tau$ at a given $x$ differ by at most $4$, so every slot of a level $j\ge q_1+3$ is strictly more expensive than every slot of the levels $j\le q_1$, and mass placed on it can be moved to a cheaper slot with spare capacity. Hence $J(x)$ is the minimum of a linear program with finitely many slots, and both minima are attained by filling the cheapest slots first (\cref{lem:lp-lipschitz}).

\begin{lemma}\label{lem:profile-regular}
Let $0<\delta<\kappa$. On $[\delta,\kappa]$ the function $\mathsf B$ is bounded, and it is continuous outside the finite set of points $s/j$ with $s\in\{a,\kappa,1,\kappa-a\}$ and $j\in\Z_{\ge1}$. On each interval between consecutive such points the weights are constant and the capacities are affine functions of $x$.
\end{lemma}

\begin{proof}
Between consecutive points of the set the integers $q_a,q_\kappa,q_1$ are constant, and so is the order of $r_a$ and $r_\kappa$, since $r_\kappa-r_a=(\kappa-a)-(q_\kappa-q_a)x$ vanishes only at $x=(\kappa-a)/(q_\kappa-q_a)$. Hence the types and weights are constant and the lengths $\ell_\tau$ are affine; the finite set of jet slots with $j\le q_1+2$ is constant as well. The minimum of a linear program with finitely many slots and fixed costs is Lipschitz in the capacities (\cref{lem:lp-lipschitz}), so $J$ and $R$ are continuous there; $\mathsf N$ is affine there. Boundedness follows from $q_\kappa\le\kappa/\delta$.
\end{proof}

\subsection{Uniform approximation}\label{sec:uniform}

\begin{proposition}\label{prop:profile-uniform}
Let $0<\delta<\kappa$, and put $C_\delta=\kappa/\delta$. For every $n$ with $h>\kappa/\delta^2$ and every ordinary prime $p$ with $\delta h\le p\le K$,
\[
\vg p(\Qt_n)\ \ge\ h\,\mathsf B(p/h)-C_\delta .
\]
More precisely $\cJ_p=h\,J(p/h)$ and $\cP_p\ge h\,R(p/h)-C_\delta$. We write $n_\delta$ for the least $n$ with $h>\kappa/\delta^2$.
\end{proposition}

\begin{proof}
Since $h>\kappa/\delta^2$ we have $p^2\ge\delta^2h^2>\kappa h=K$, so \cref{prop:jet,prop:pole} apply. Put $x=p/h\in[\delta,\kappa]$, and note that $a/x=A/p$ and $\kappa/x=K/p$.

\emph{Normalization.} Since $h\le K<p^2$, Legendre's formula gives $v_p(m!)=\lfloor m/p\rfloor$ for $m\le K$. Hence $v_p(S_n)=\lfloor K/p\rfloor-3\lfloor A/p\rfloor=q_\kappa-3q_a$, and $\sum_{i<h}v_p(i!)=\sum_{i<h}\lfloor i/p\rfloor=p\,q_1(q_1-1)/2+q_1(h-q_1p)$, with $q_1=\lfloor h/p\rfloor$. Therefore $v_p(S_n^h/F_n)=h\,\mathsf N(x)$ exactly.

\emph{Class counts.} The numbers $hr_a=A-q_ap$ and $hr_\kappa=K-q_\kappa p$ are integers in $[0,p-1]$. For $1\le\sigma<p$ we have $a_\sigma=q_a+\ind_{\sigma\le hr_a}$ and $c_\sigma=q_\kappa+\ind_{\sigma\le hr_\kappa}$, and for the zero class $a_0=q_a$ and $c_0=q_\kappa$. Cut the interval $(0,x]$ into the $p$ intervals $\mathcal I_\sigma=((\sigma-1)/h,\sigma/h]$, $1\le\sigma\le p$. Since $r_a$ and $r_\kappa$ are multiples of $1/h$, the pair $(\alpha,\gamma)$ is constant on the interior of each $\mathcal I_\sigma$: for $\sigma<p$ it equals $(a_\sigma,c_\sigma)$, and on $\mathcal I_p$, which lies to the right of $r_a$ and $r_\kappa$, it equals $(q_a,q_\kappa)=(a_0,c_0)$. Say that a class $\sigma$ has type $\tau$ if $(a_\sigma,c_\sigma)=\tau$, and let $N_\tau$ be the number of classes $0\le\sigma<p$ of type $\tau$. Matching the class $\sigma\ge1$ with $\mathcal I_\sigma$ and the zero class with $\mathcal I_p$, we obtain
\begin{equation}\label{eq:class-counts}
N_\tau=h\,\ell_\tau(x)\qquad\text{for every type }\tau,
\end{equation}
with no error term. In particular the zero class has the type $\tau_0=(q_a,q_\kappa)$, which occurs at $x$.

\emph{The jet bound.} A class of type $\tau$ has $w_\sigma=\omega_\tau$. By \eqref{eq:class-counts} the multiset $\cS^{\mathrm{jet}}_p$ contains the weight $\omega_\tau+2j$ of the jet slot $(\tau,j)$ exactly $h\ell_\tau(x)$ times, and nothing else. The sum of its $h$ smallest elements is therefore
\[
\cJ_p=\min\Bigl\{\textstyle\sum_{\tau,j}(\omega_\tau+2j)\,Y_{\tau,j}\ :\ Y_{\tau,j}\in\Z,\ 0\le Y_{\tau,j}\le h\ell_\tau(x),\ \sum_{\tau,j}Y_{\tau,j}=h\Bigr\}.
\]
Each class contributes exactly one element to each level $j$, so the number of elements of level $j$ and type $\tau$ is $N_\tau=h\ell_\tau(x)$; the discrete capacities are thus the exact capacities of \cref{def:profile} multiplied by $h$, with no perturbation. The argument after \cref{def:profile} therefore applies verbatim to the discrete program (the levels $j\le q_1$ hold $(q_1+1)p>h$ elements and the weights at a given $p$ differ by at most $4$), so only the finitely many slots with $j\le q_1+2$ matter, here and in the program defining $J(x)$. The same holds for the zero class, whose type $\tau_0$ occurs at $x$ by \eqref{eq:class-counts}, whatever $a_0$ and $c_0$ are; in particular no estimate of the difference between the class counts and the capacities is needed, and the constant $C_\delta=\kappa/\delta$ of \cref{prop:profile-uniform} depends only on $\delta$ and the triple $(H,B,C)$, which is what \cref{thm:arith-general} uses. For this finite program the minimum over real $Y_{\tau,j}$ is attained by filling the slots in order of increasing weight (\cref{lem:lp-lipschitz}); as the capacities $h\ell_\tau(x)$ and the total $h$ are integers, this optimal solution is integral, so the integrality constraint can be dropped. The substitution $Y=hy$ then turns the program into $h$ times the program defining $J(x)$, and $\cJ_p=hJ(x)$.

\emph{The pole bound.} For a finite multiset $\cS$ of integers, $\Sigma^-_h(\cS)$ is the minimum of the sums of the sub-multisets of $\cS$ with at most $h$ elements. Let $\cS'$ be obtained from $\cS^{\mathrm{pole}}_p$ by replacing the $m_0=c_0-a_0=q_\kappa-q_a$ elements $w_0+1+2j$ of the zero class by the elements $w_0+2j+2\cdot\ind_{a_0=0,\ j=m_0-1}$ ($0\le j<m_0$) that a nonzero class of the same type $\tau_0$ would contribute. Each of these elements changes by $+1$ or $-1$, so $\Sigma^-_h(\cS^{\mathrm{pole}}_p)\ge\Sigma^-_h(\cS')-m_0$. By \eqref{eq:class-counts}, $\cS'$ contains the weight of the pole slot $(\tau,j)$ exactly $h\ell_\tau(x)$ times, and the argument given for the jet bound, with $\sum Y\le h$ in place of $\sum Y=h$, shows $\Sigma^-_h(\cS')=hR(x)$. Since $m_0\le q_\kappa\le\kappa/x\le\kappa/\delta$, we get $\cP_p\ge hR(x)-C_\delta$.

\emph{Conclusion.} Combining with \cref{prop:jet,prop:pole},
\[
\vg p(\Qt_n)=v_p(S_n^h/F_n)+\vg p(Q_n)\ge h\mathsf N(x)+\max(\cJ_p,\cP_p)\ge h\mathsf B(x)-C_\delta. \qedhere
\]
\end{proof}

The identity $\cJ_p=hJ(p/h)$ and the inequality $|\cP_p-hR(p/h)|\le q_\kappa-q_a$ were also checked directly, from the definitions of the class counts and of \cref{def:profile}, for the four triples, for $n\le30$ and $n\in\{47,64,101\}$, and for all ordinary primes $\sqrt K<p\le K$ ($3875$ cases, exact rational arithmetic).

\subsection{Summation over primes}\label{sec:pnt}

\begin{lemma}\label{lem:pnt}
Let $0<\delta<\kappa$ and let $\phi$ be a bounded function on $[\delta,\kappa]$, continuous outside a finite set. Then
\[
\lim_{h\to\infty}\frac1h\sum_{\delta h\le p\le\kappa h}\phi(p/h)\log p=\int_\delta^\kappa\phi(x)\dd x .
\]
\end{lemma}

\begin{proof}
By the prime number theorem, $\vartheta(y)=\sum_{p\le y}\log p=y+o(y)$. For the indicator of an interval $[\alpha,\beta]\subset[\delta,\kappa]$ the sum equals $(\vartheta(\beta h)-\vartheta(\alpha h)+O(\log h))/h\to\beta-\alpha$, so the claim holds for step functions. A function as in the statement is Riemann integrable, hence squeezed between step functions with integrals arbitrarily close to its own; since the weights $\log p/h$ are nonnegative, the claim follows.
\end{proof}

\begin{theorem}[arithmetic estimate, general form]\label{thm:arith-general}
For every $\delta\in(0,\kappa)$,
\[
\liminf_{n\to\infty}\frac1{h^2}\sum_p\vg p(\Qt_n)\log p\ \ge\ I_\delta+\eps_{\mathrm{exc}},\qquad I_\delta=\int_\delta^\kappa\mathsf B(x)\dd x-2\delta,
\]
with $\eps_{\mathrm{exc}}$ as in \cref{cor:exceptional}. For the indices with $\Qt_n\ne0$ the sum is $\log\cont(\Qt_n)$ by \eqref{eq:content-sum}; for the other indices every $\vg p(\Qt_n)$ is $+\infty$. Hence the lower limit of $h^{-2}\log\cont(\Qt_n)$ along the indices with $\Qt_n\ne0$ is at least $I_\delta+\eps_{\mathrm{exc}}$.
\end{theorem}

\begin{proof}
Fix $\delta$ and let $n\ge n_\delta$ be large; in particular $\delta h>\max(\Ex)$ and $\delta h>\sqrt N$. The bounds below hold for every such $n$, with $\vg p(\Qt_n)=+\infty$ allowed. Split the primes into four sets.
\begin{enumerate}[(i)]
\item $p\in\Ex$: by \cref{cor:exceptional} these contribute at least $\eps_{\mathrm{exc}}h^2-O(h\log h)$.
\item Ordinary $p<\delta h$: by \cref{prop:generic} the contribution is at least $-2h\sum_{p<\delta h}L_p\log p$. For $p\le\sqrt N$ we use $L_p\log p\le\log N$, and for $\sqrt N<p<\delta h$ we have $L_p=1$. The total is at least $-2h\bigl(\pi(\sqrt N)\log N+\vartheta(\delta h)\bigr)=-2\delta h^2-o(h^2)$.
\item Ordinary $p$ with $\delta h\le p\le K$: by \cref{prop:profile-uniform,lem:profile-regular,lem:pnt} the contribution is at least $h\sum_{\delta h\le p\le\kappa h}\mathsf B(p/h)\log p-C_\delta\vartheta(K)=h^2\int_\delta^\kappa\mathsf B-o(h^2)$. Here every prime of the range is ordinary because $\delta h>\max(\Ex)$, and $C_\delta\vartheta(K)=O(h)$ because $C_\delta=\kappa/\delta$ does not depend on $p$ or $n$.
\item Ordinary $p>K$: the contribution is nonnegative by \cref{prop:generic}.
\end{enumerate}
Adding the four contributions and dividing by $h^2$ proves the theorem.
\end{proof}

\subsection{The certified integrals}\label{sec:profile-values}
The integral $I_\delta$ is bounded below by a finite computation. By \cref{lem:profile-regular}, on each subinterval of the partition of $[\delta,\kappa]$ by the points $s/j$ the weights are fixed and the capacities affine. The minimum of a linear program is a convex function of its capacities, so $J$ and $R$, and hence $\mathsf B$, are convex on each subinterval (\cref{lem:profile-convex}); by Jensen's inequality the midpoint rule therefore gives a lower bound for the integral over each subinterval. The certificate evaluates $\mathsf B$ exactly at the midpoints (\cref{sec:cert-profile}). The resulting rational lower bounds are recorded in \cref{tab:profile}.

\begin{proposition}[profile certificate]\label{prop:profile-certificate}
With $\delta$ as in \cref{tab:profile}, $I_\delta$ exceeds the value in the third column of that table.
\end{proposition}

\begin{proof}
By \cref{lem:profile-convex}, $I_\delta+2\delta\ge\sum_{\mathcal I}|\mathcal I|\,\mathsf B(x_{\mathcal I})$, the sum running over the subintervals of the partition of \cref{sec:cert-profile}. Each midpoint value $\mathsf B(x_{\mathcal I})$ is a rational number computed exactly from \cref{def:profile} by filling the cheapest slots first; the finite sum is evaluated in exact rational arithmetic and compared with the values of \cref{tab:profile}. The step is certified in the sense of \cref{sec:intro-status}: it was reproduced by the original programs and confirmed, with the exact value of $I_\delta$, by a separate recomputation.
\end{proof}

\begin{table}[ht]
\centering
\caption{The certified lower bounds for $I_\delta=\int_\delta^\kappa\mathsf B-2\delta$ (normalized by $h^2$), and the resulting arithmetic rates $I_\delta+\eps_{\mathrm{exc}}$. Decimals are rounded down. For $\lambda=2/11$ the certificate is computed in the scale $y=p/n=4x$, where it reads $\int_{1/40}^6\mathsf B_n(y)\dd y-\frac8{40}>\frac92$ with $\mathsf B_n(y)=4\mathsf B(y/4)$; this is the bound $I_{1/160}>9/32$.}
\label{tab:profile}
\begin{tabular}{@{}lcccc@{}}
\toprule
$\lambda$ & $\delta$ & lower bound for $I_\delta$ & $\eps_{\mathrm{exc}}$ & $I_\delta+\eps_{\mathrm{exc}}$\\
\midrule
$2/11$ & $1/160$ & $0.2816209$ & $-0.5523150$ & $-0.2706941$\\
$2/9$ & $1/500$ & $-0.8850553$ & $-0.4394723$ & $-1.3245276$\\
$-2/9$ & $1/500$ & $0.1673093$ & $-0.5503976$ & $-0.3830883$\\
$-2/7$ & $1/500$ & $0.0769450$ & $-0.5982677$ & $-0.5213227$\\
\bottomrule
\end{tabular}
\end{table}

\begin{proof}[Proof of \cref{thm:arith}]
By \cref{thm:arith-general,prop:profile-certificate}, $\liminf h^{-2}\log\cont(\Qt_n)\ge I_\delta+\eps_{\mathrm{exc}}$, and the certificates verify, with outward-rounded rational enclosures of the logarithms in $\eps_{\mathrm{exc}}$, that $I_\delta+\eps_{\mathrm{exc}}>\rar(\lambda)$ for the thresholds of \cref{tab:rates}. (For $\lambda=2/11$ the certificate uses the weaker bound $I_\delta>9/32$, which gives $-0.271065>-2169/8000$.)
\end{proof}

\begin{remark}\label{rem:profile-float}
The AI-internal recomputation described in the section on reproducibility evaluates $I_\delta$ exactly, as a rational number: in the $n$-scale of $2/11$, $\int_{1/40}^6\mathsf B_n-\frac15=4.5059344662496163\ldots$, and $I_\delta=-0.8850552343712680\ldots$, $0.1685590216\ldots$ and $0.0781087782\ldots$ for $2/9$, $-2/9$ and $-2/7$. For $2/11$ and $2/9$ the function $\mathsf B$ is affine on every subinterval, so the midpoint values of \cref{tab:profile} are exact there; for $-2/9$ and $-2/7$ the midpoint values are lower bounds, by convexity, and they are conservative by about $0.0012$. A midpoint rule whose nodes are not aligned with the partition of \cref{sec:cert-profile} is not a one-sided bound and is not used.
\end{remark}

\section{Nonvanishing at prime indices}\label{sec:nonvanishing}

This section proves \cref{thm:nonvanishing}, the third hypothesis of the criterion \cref{lem:criterion}. In that criterion the determinant is evaluated at a fixed rational number, the hypothetical value of $\xi$, so it suffices to show that $\Qt_n(r)\ne0$ for each fixed $r\in\Q$ and infinitely many $n$. The reduction (\cref{thm:nonvanishing-general}) is a candidate and the local determinants (\cref{prop:local-dets}) are certified, in the sense of \cref{sec:intro-status}. We take $n=p$ prime and use the residue classes modulo the same prime $p$. Then $A=Bp$ and $K=Cp$, so that every class contains exactly $B$ zeros and $C$ poles of $W_p$, and the Gram matrix decomposes into $p$ nearly independent blocks. The reduction of nonvanishing to local Hankel determinants of a local functional, through residue classes modulo a prime and a basis adapted to them, is the method of \cite[Section~5]{Litsey2026Li2}, where the prime runs through a band depending on $n$ and the local determinants are called type constants; taking the index itself prime leaves only the two determinants $\det\mathsf H_0$, $\det\mathsf H_1$.

\subsection{The local functional}\label{sec:local-functional}
Let $\cV$ be the $\Q$-linear functional on rational functions in $u$ with at most simple poles in $\{0,-1,-2,\dots\}$ defined by
\[
\cV(u^k)=k\,M_{k-1}(\lambda)\quad(k\ge1),\qquad \cV(1)=0,\qquad \cV\Bigl(\frac1{u+k}\Bigr)=\lambda^{-k}\tau_k(\lambda)\quad(k\ge0).
\]
Equivalently, $\cV(\rho+\sum_kr_k(u+k)^{-1})=G(\rho')+\sum_kr_k\lambda^{-k}\tau_k(\lambda)$ for a polynomial $\rho$; note $\tau_0=0$. Define the local weights and the local Hankel matrices
\[
w_1(u)=\frac{\prod_{k=0}^{B-1}(u+k)^2}{\prod_{k=B}^{C-1}(u+k)},\qquad w_0(u)=\frac{u\prod_{k=1}^{B}(u+k)^2}{\prod_{k=B+1}^{C}(u+k)},
\]
\[
\mathsf H_b=\bigl[\cV(u^{i+j}w_b)\bigr]_{0\le i,j<H}\qquad(b=0,1).
\]
For example, for $\lambda=2/11$ one has $w_1=u^2/\prod_{k=1}^5(u+k)$ and $w_0=u(u+1)^2/\prod_{k=2}^6(u+k)$, and for $\lambda=2/9$ one has $w_1=1/\prod_{k=0}^8(u+k)$ and $w_0=u/\prod_{k=1}^9(u+k)$.

\begin{proposition}[local determinants]\label{prop:local-dets}
The determinants $\det\mathsf H_1$ and $\det\mathsf H_0$ are the nonzero rational numbers listed in \cref{tab:local-dets}.
\end{proposition}

\begin{proof}
This is a finite computation in exact rational arithmetic. The moments $M_k(\lambda)$ are rational by \eqref{eq:moments}, the partial fractions of $u^{i+j}w_b$ are obtained by polynomial division and residues, and the determinants by exact elimination; each determinant is confirmed by an explicit factorization $\mathsf H_b=\mathsf L\mathsf U$ checked entrywise, and all eight values were recomputed by a separate program (\cref{sec:cert-local}). The step is certified in the sense of \cref{sec:intro-status}.
\end{proof}

\begin{table}[ht]
\centering
\caption{The local determinants.}
\label{tab:local-dets}
\small
\begin{tabular}{@{}llc@{}}
\toprule
$\lambda$ & & value\\
\midrule
$2/11$ & $\det\mathsf H_1$ & $\dfrac{17828022563089333969}{1224440064}$\\[9pt]
 & $\det\mathsf H_0$ & $\dfrac{12598972834814973210048467}{2295825120}$\\[9pt]
$2/9$ & $\det\mathsf H_1$ & $\dfrac{3951366107675153857566158184061855291890957151}{50404062960731160576000000000000}$\\[9pt]
 & $\det\mathsf H_0$ & $\dfrac{75748065809620696385810762379349631101207641762772048941}{23344937339644196864000000000000}$\\[9pt]
$-2/9$ & $\det\mathsf H_1$ & $\dfrac{1677163510311484581}{15488000}$\\[9pt]
 & $\det\mathsf H_0$ & $\dfrac{16417954848664067036722569}{524733440}$\\[9pt]
$-2/7$ & $\det\mathsf H_1$ & $\dfrac{14067971871496713969085291850153}{1057916215296000000}$\\[9pt]
 & $\det\mathsf H_0$ & $\dfrac{17800469908791563702231968746821382359}{495898225920000000}$\\[9pt]
\bottomrule
\end{tabular}
\end{table}

\subsection{The reduction}\label{sec:reduction}

\begin{theorem}\label{thm:nonvanishing-general}
Let $r\in\Q$, and let $p$ be a prime with $p>C$, $p\notin\Ex$, $p$ not dividing the denominator of $r$, and $p$ not dividing the numerators of $\det\mathsf H_0$ and $\det\mathsf H_1$. Then $Q_p(r)\ne0$.
\end{theorem}

The hypotheses are used as follows. The condition $p>C$ gives $K=Cp<p^2$, the standing assumption of \cref{sec:ordinary}, and makes the local poles $k\le C$ of the two local weights distinct modulo $p$. The condition $p\notin\Ex$ makes $p$ odd and $\lambda$, $1-\lambda$ units. Under these two conditions the entries of $\mathsf H_0$ and $\mathsf H_1$ lie in $\Zp$: the polynomial part of $u^{i+j}w_b$ has integer coefficients, its residues are integers divided by products of differences of poles, and $\tau_k(\lambda)\in\Zp$ for $k\le C<p$. So $\det\mathsf H_b\in\Zp$, and no condition on its denominator is needed. The condition on $r$ allows the specialization $X=r$ in $\Zp$.

\begin{proof}
Let $n=p$, so $h=Hp$, $A=Bp$, $K=Cp<p^2$ because $p>C$, and for every class $a_\sigma=B$, $c_\sigma=C$, $w_\sigma=3B-C-1=:w$.

\emph{Step 1: an adapted basis.} For $0\le\sigma<p$ and $0\le i<H$ put
\[
e_{\sigma,i}(x)=(x+\sigma)^i\prod_{\rho\ne\sigma}(x+\rho)^H,
\]
the product running over $0\le\rho<p$. These are $Hp=h$ polynomials of degree less than $h$ with integer coefficients. Modulo $p$, the ideals $((x+\rho)^H)$ are pairwise coprime and $\prod_\rho(x+\rho)^H\equiv(x^p-x)^H$ has degree $h$, so by the Chinese remainder theorem the reductions of polynomials of degree $<h$ correspond bijectively to $\prod_\rho\mathbb F_p[x]/((x+\rho)^H)$. Under this correspondence $e_{\sigma,i}$ has zero components except at $\sigma$, where it is $(x+\sigma)^i$ times a unit. Hence the $e_{\sigma,i}$ reduce to a basis, and the matrix $E\in\Z^{h\times h}$ expressing them in the monomial basis has $\det E\in\Zp^\times$. Consequently
\[
\mathsf G:=\bigl[\Lambda_r(W_p\,e_{\sigma,i}e_{\tau,j})\bigr]_{(\sigma,i),(\tau,j)}=E\,[\Lambda_r(W_px^{i+j})]\,E^{\mathsf T},\qquad \det\mathsf G=(\det E)^2Q_p(r).
\]

\emph{Step 2: class decomposition.} Fix $\mu$ large (as specified below). By \cref{prop:distribution} and specialization $X=r\in\Zp$, each entry of $\mathsf G$ is congruent modulo $p^{\mu-c}$ to $\sum_{\rho}\Theta_\rho\bigl(e_{\sigma,i}(pu-\rho)e_{\tau,j}(pu-\rho)\bigr)$, with $c$ independent of $\mu$. Now $e_{\sigma,i}(pu-\rho)$ has minimal coefficient valuation exactly $i$ if $\rho=\sigma$ (it is $p^iu^i$ times a polynomial with unit constant term), and at least $H$ if $\rho\ne\sigma$ (it contains the factor $(pu)^H$). By \cref{lem:local-integrality} the contribution of the class $\rho$ to the entry $((\sigma,i),(\tau,j))$ has valuation at least
\[
w+\ind_{\rho=0}+\bigl(i\ \text{if}\ \rho=\sigma,\ \text{else}\ H\bigr)+\bigl(j\ \text{if}\ \rho=\tau,\ \text{else}\ H\bigr).
\]

\emph{Step 3: scaling.} Let $\mathsf Z$ be obtained from $\mathsf G$ by multiplying the entry $((\sigma,i),(\tau,j))$ by $p^{-w-i-j}$ and then dividing the $H$ rows $(0,i)$ by $p$. Then $\det\mathsf Z=p^{-e_0}\det\mathsf G$ for an explicit integer $e_0$. By Step 2, the contribution of the class $\rho$ to $\mathsf Z_{(\sigma,i),(\tau,j)}$ has valuation at least $\ind_{\rho=0}-\ind_{\sigma=0}$ plus $0$, $H-j$, $H-i$ or $2H-i-j$ according as $\rho=\sigma=\tau$, $\rho=\sigma\ne\tau$, $\rho=\tau\ne\sigma$ or $\rho\notin\{\sigma,\tau\}$. Since $i,j\le H-1$, it follows that every entry of $\mathsf Z$ lies in $\Zp$ (modulo $p^{\mu-c-w'}$ with $w'=|w|+2H+1$, which is harmless for $\mu$ large), that the entries in rows $(\sigma,i)$ with $\sigma\ne0$ and columns $(\tau,j)$ with $\tau\ne\sigma$ are divisible by $p$, and that in a diagonal block $\sigma=\tau$ only the own class $\rho=\sigma$ contributes modulo $p$. Thus, modulo $p$, $\mathsf Z$ is block upper triangular when the zero class is listed first, and
\[
\det\mathsf Z\equiv\prod_{\sigma=0}^{p-1}\det\mathsf Z^{(\sigma)}\pmod p,\qquad \mathsf Z^{(\sigma)}_{ij}\equiv p^{-w-i-j-\ind_{\sigma=0}}\ \Theta_\sigma\bigl(e_{\sigma,i}(pu-\sigma)e_{\sigma,j}(pu-\sigma)\bigr).
\]

\emph{Step 4: the diagonal blocks.} Collect the factors as in the proof of \cref{lem:local-integrality}. For $\sigma\ne0$ the local zeros and poles are $p(u+k)$ with $0\le k<B$, respectively $0\le k<C$; for $\sigma=0$ they are $p(u+k)$ with $1\le k\le B$, respectively $1\le k\le C$, and the factor $x=pu$ contributes $pu$. All other factors, including those of $e_{\sigma,i}(pu-\sigma)e_{\sigma,j}(pu-\sigma)$ outside $(pu)^{i+j}$ and, for $\sigma\ne0$, the factor $pu-\sigma$, are polynomials in $u$ with unit constant term and all other coefficients divisible by $p$. Hence
\[
(pu-\sigma)W_p(pu-\sigma)\,e_{\sigma,i}(pu-\sigma)e_{\sigma,j}(pu-\sigma)=p^{\,w+1+\ind_{\sigma=0}+i+j}\,k_\sigma\,u^{i+j}w_{b}(u)\bigl(1+p\,g_{\sigma}(u)\bigr),
\]
with $b=\ind_{\sigma\ne0}$, a unit $k_\sigma\in\Zp^\times$, and $g_{\sigma}=g_1/D_1$ where $g_1,D_1\in\Zp[u]$ and $D_1\equiv D_1(0)\not\equiv0\pmod p$. The function $u^{i+j}w_bg_\sigma$ is of the form $\Phi_1/(BD)$ considered in the proof of \cref{lem:local-integrality}, with $\Phi_1\in\Zp[u]$; the argument given there shows that $T_{M,\lambda_p,Y_\mu}$ maps its derivative into $\Zp[X]$, so its contribution to $\mathsf Z^{(\sigma)}_{ij}$ is divisible by $p$. Therefore
\[
\mathsf Z^{(\sigma)}_{ij}\equiv\lambda^{-\sigma}k_\sigma\,T_{M,\lambda_p,Y_\mu}\Bigl(\frac{\mathrm d}{\mathrm du}\bigl[u^{i+j}w_b(u)\bigr]\Bigr)\pmod p .
\]
Write $u^{i+j}w_b=\rho+\sum_k\varrho_k(u+k)^{-1}$, where the $k$ are the local poles ($k\le C<p$) and $\rho\in\Zp[u]$, $\varrho_k\in\Zp$ because the differences of the poles are units. By \cref{lem:finite-moments,lem:pole-regularization},
\[
T_{M,\lambda_p,Y_\mu}\bigl(\rho'-\textstyle\sum_k\varrho_k(u+k)^{-2}\bigr)\equiv G_{\lambda_p}(\rho')-\sum_k\varrho_k\lambda_p^{-k}\bigl(Y_\mu-\tau_k(\lambda_p)\bigr)\pmod{p^\mu}.
\]
Since $Y_\mu\equiv0\pmod{p^2}$ at $X=r$, $\lambda_p\equiv\lambda\pmod p$, and $G_{\lambda_p}(\rho')$, $\tau_k(\lambda_p)$ are rational functions of $\lambda_p$ whose denominators are units, the right side is congruent modulo $p$ to $G(\rho')+\sum_k\varrho_k\lambda^{-k}\tau_k(\lambda)=\cV(u^{i+j}w_b)$. Hence $\det\mathsf Z^{(\sigma)}\equiv(\lambda^{-\sigma}k_\sigma)^H\det\mathsf H_b\pmod p$.

\emph{Step 5: conclusion.} Choose $\mu$ so large that all congruences above hold modulo $p$ for the exact matrix (the scaling exponents are bounded by $|w|+2H+1$, independently of $\mu$). The numbers $\det\mathsf H_0$ and $\det\mathsf H_1$ lie in $\Zp$, as noted after the statement, and $p$ does not divide their numerators; so they are $p$-adic units. Therefore $\det\mathsf Z\not\equiv0\pmod p$, so $\det\mathsf G\ne0$ and $Q_p(r)\ne0$.
\end{proof}

\begin{proof}[Proof of \cref{thm:nonvanishing}]
By \cref{prop:local-dets} the numbers $\det\mathsf H_0$ and $\det\mathsf H_1$ are nonzero rationals, so all but finitely many primes satisfy the hypotheses of \cref{thm:nonvanishing-general}. For those primes $Q_p(r)\ne0$, and $\Qt_p(r)=S_p^h F_p^{-1}Q_p(r)\ne0$.
\end{proof}

\begin{remark}
Only nonvanishing at a fixed rational point is needed: the final argument evaluates $\Qt_p$ at the hypothetical rational value of $\xi$, not at the real number $\xi$ itself. Nonvanishing of $Q_p$ as a polynomial is a consequence. The proof gives more: under the hypotheses of \cref{thm:nonvanishing-general}, $v_p(Q_p(r))=e_0$ exactly, where $e_0=hw+pH(H-1)+H$ is the exponent of Step~3. On the exact polynomials this equality was observed at the primes satisfying the hypotheses that were tested ($p=7,13,19$ for $2/11$, $p=11$ for $2/9$, $p=13,17,19$ for $-2/9$, $p=13,17$ for $-2/7$, at $r=0,1,-3/5$), and it fails at $p=17$ for $2/11$ and at $p=11$ for $-2/7$, where the valuation is $e_0+1$ and $e_0+10$; these two primes divide the numerator of $\det\mathsf H_0$ and of $\det\mathsf H_1$ respectively, and are excluded by the hypothesis. These observations are numerical evidence and not part of the proof.
\end{remark}

\section{Archimedean decay I: logarithmic energy and positive arguments}\label{sec:energy}

This section and the next bound $\log|\Qt_n(\xi)|$ from above in terms of a comparison measure; with the measures of \cref{app:certificates} this gives \cref{thm:arch} (\cref{sec:proof-arch}). The two ingredients common to both are a regularized energy inequality (\cref{sec:energy-inequality}) and Stirling's formula for the normalization (\cref{lem:stirling}). For positive $\lambda$ we then use Heine's formula for the derivative functional of \cref{lem:series} directly (\cref{prop:positive-bound}). The statements are candidates in the sense of \cref{sec:intro-status}. The structure of the argument is that of \cite[Section~6]{Litsey2026Li2}: Heine's formula for a functional with derivatives, in which a site carries at most two variables (Lemma~6.2 there), a comparison with the energy of a measure, and a bound with a level $\theta$ and a penalty integral (Theorem~6.1 there, for the lattice). The energy inequality is proved here by a different, elementary route (\cref{lem:neutral-energy,lem:energy-comparison}).

\subsection{An energy inequality}\label{sec:energy-inequality}
For $\eps>0$ and $z,w\in\C$ put $\mathcal K_\eps(z,w)=\tfrac12\log(|z-w|^2+\eps^2)$.

\begin{lemma}[neutral charges]\label{lem:neutral-energy}
For $\eps>0$, points $z_1,\dots,z_m\in\C$ and real numbers $b_1,\dots,b_m$ with $\sum_ib_i=0$,
\[
\sum_{i,j}b_ib_j\,\mathcal K_\eps(z_i,z_j)\le0 .
\]
\end{lemma}

\begin{proof}
For $\theta>0$ the matrix $[e^{-\theta|z_i-z_j|^2}]$ is positive semidefinite: writing $z=x+\mathrm iy$, it equals the congruence by the diagonal matrix $\mathrm{diag}(e^{-\theta|z_i|^2})$ of the matrix $[e^{2\theta(x_ix_j+y_iy_j)}]$, which is a convergent sum of the Hadamard powers $[(x_ix_j+y_iy_j)^k]$ with nonnegative coefficients, each positive semidefinite by Schur's product theorem. Since $(s+\varrho^2)^{-1}=\int_0^\infty e^{-\theta s}e^{-\theta\varrho^2}\dd\theta$ for $s>0$, the matrix $[(s+|z_i-z_j|^2)^{-1}]$ is positive semidefinite. Hence $\varphi(s)=\sum_{i,j}b_ib_j\log(s+|z_i-z_j|^2)$ is nondecreasing on $(0,\infty)$. As $\sum b_i=0$, $\varphi(s)=\sum_{i,j}b_ib_j\log(1+|z_i-z_j|^2/s)\to0$ when $s\to\infty$. So $\varphi(\eps^2)\le0$.
\end{proof}

A \emph{comparison measure} is a positive Borel measure $\sigma$ on $\C$ of total mass $1$, supported on a compact subset of a finite union of $m_\sigma$ straight lines, and absolutely continuous with respect to arc length on each line with density at most $\varrho_\sigma$. Its potential and energy are
\[
U_\sigma(z)=\int\log|z-w|\dd\sigma(w),\qquad E(\sigma)=\int U_\sigma\dd\sigma .
\]
The potential is continuous on $\C$, and $E(\sigma)$ is finite.

\begin{lemma}[energy comparison]\label{lem:energy-comparison}
Let $\sigma$ be a comparison measure. For every $h\ge2$ and all $z_1,\dots,z_h\in\C$,
\[
\sum_{\substack{i\ne j\\ z_i\ne z_j}}\log|z_i-z_j|\ \le\ -h^2E(\sigma)+2h\sum_{i=1}^hU_\sigma(z_i)+2\,n_c\log h+2\pi m_\sigma\varrho_\sigma,
\]
where $n_c$ is the number of ordered pairs $(i,j)$, including $i=j$, with $z_i=z_j$.
\end{lemma}

\begin{proof}
Let $\eps=h^{-2}$, $U^\eps(z)=\int \mathcal K_\eps(z,w)\dd\sigma(w)$ and $E^\eps=\int U^\eps\dd\sigma$. Approximate $\sigma$ weakly by finite positive point measures $\sigma_k$ of mass one supported in $\operatorname{supp}\sigma$ (Riemann sums on each line). Since $\mathcal K_\eps$ is continuous, \cref{lem:neutral-energy} applied to $\mu-\sigma_k$, with $\mu=h^{-1}\sum_i\delta_{z_i}$, gives in the limit $\iint \mathcal K_\eps\dd(\mu-\sigma)\dd(\mu-\sigma)\le0$, that is
\[
\sum_{i,j}\mathcal K_\eps(z_i,z_j)\le2h\sum_iU^\eps(z_i)-h^2E^\eps .
\]
On the left, $\mathcal K_\eps(z,w)\ge\log|z-w|$ for $z\ne w$ and $\mathcal K_\eps(z,z)=\log\eps=-2\log h$. On the right, $E^\eps\ge E(\sigma)$, and $0\le U^\eps(z)-U_\sigma(z)\le\pi m_\sigma\varrho_\sigma\eps$: on each line, parametrized by arc length $u$, with $u_0$ the parameter of the orthogonal projection of $z$, we have $|z-w|\ge|u-u_0|$ and $\int_\R\frac12\log\bigl(1+\eps^2(u-u_0)^{-2}\bigr)\dd u=\pi\eps$. Combining these estimates proves the lemma.
\end{proof}

\begin{lemma}[normalization]\label{lem:stirling}
As $n\to\infty$,
\[
\log\frac{S_n^h}{F_n}=(d-1)h^2\log h+\Cst h^2+O(h\log h),\qquad \Cst=\kappa\log\kappa-3a\log a+\tfrac32-d,
\]
with $0\log0=0$.
\end{lemma}

\begin{proof}
By Stirling's formula $\log m!=m\log m-m+O(\log(m+1))$, and $\sum_{i<h}(i\log i-i)=\frac{h^2}2\log h-\frac34h^2+O(h\log h)$. Hence $h\log K!-3h\log A!=(\kappa-3a)h^2\log h+(\kappa\log\kappa-\kappa-3a\log a+3a)h^2+O(h\log h)$ and $\log F_n=h^2\log h-\frac32h^2+O(h\log h)$.
\end{proof}

The constants are $\Cst=2.3979184$, $0.5075059$, $2.5190484$ and $2.2176686$ for $\lambda=2/11$, $2/9$, $-2/9$ and $-2/7$.

\subsection{Heine's formula with derivatives}\label{sec:heine}
For a linear functional $\ell$ on $\C[x]$ let $\ell^{\otimes h}$ be the linear functional on $\C[x_1,\dots,x_h]$ with $\ell^{\otimes h}(\prod_ix_i^{k_i})=\prod_i\ell(x^{k_i})$, and let $\Delta(x_1,\dots,x_h)=\prod_{i<j}(x_j-x_i)$.

\begin{lemma}[Heine]\label{lem:heine}
$\det_{0\le i,j<h}\ell(x^{i+j})=\frac1{h!}\,\ell^{\otimes h}(\Delta^2)$.
\end{lemma}

\begin{proof}
Expanding $\Delta=\det[x_i^{j}]$ over permutations,
\[
\Delta=\sum_\pi\operatorname{sgn}(\pi)\prod_ix_i^{\pi(i)},\qquad \ell^{\otimes h}(\Delta^2)=\sum_{\pi,\pi'}\operatorname{sgn}(\pi\pi')\prod_i\ell\bigl(x^{\pi(i)+\pi'(i)}\bigr),
\]
and for each fixed $\pi$ the sum over $\pi'$ equals $\det[\ell(x^{i+j})]$ after reordering the rows.
\end{proof}

By \cref{lem:series}, the functional $\ell(\phi)=\Lambda_\xi(W_n\phi)$ on polynomials is
\[
\ell(\phi)=\sum_{x\ge1}\bigl(\alpha_x\phi(x)+\beta_x\phi'(x)\bigr),\qquad \alpha_x=\lambda^x\bigl(xW_n(x)\bigr)',\quad \beta_x=\lambda^xxW_n(x),
\]
since $(xW_n\phi)'=(xW_n)'\phi+xW_n\phi'$. The series converges absolutely for every polynomial $\phi$, so that $\ell^{\otimes h}=\sum_{x_1,\dots,x_h\ge1}\ell_{x_1}\otimes\dots\otimes\ell_{x_h}$ on polynomials, with $\ell_x(\phi)=\alpha_x\phi(x)+\beta_x\phi'(x)$, and
\begin{equation}\label{eq:heine-sum}
Q_n(\xi)=\frac1{h!}\sum_{\mathbf x\in\Z_{\ge1}^h}\mathfrak T(\mathbf x),\qquad \mathfrak T(\mathbf x)=(\ell_{x_1}\otimes\dots\otimes\ell_{x_h})(\Delta^2).
\end{equation}

\begin{lemma}[multiplicity]\label{lem:multiplicity}
There is a constant $C_0$, depending only on $\lambda$ and $(H,B,C)$, with the following property. If some value occurs at least three times among $x_1,\dots,x_h$, then $\mathfrak T(\mathbf x)=0$. Otherwise
\[
|\mathfrak T(\mathbf x)|\ \le\ (C_0h)^h\prod_{i=1}^h|\beta_{x_i}|\prod_{\substack{i\ne j\\ x_i\ne x_j}}|x_i-x_j| .
\]
\end{lemma}

\begin{proof}
For a site $s$ let $I_s=\{i:x_i=s\}$ and $m_s=|I_s|$. Write $\Delta^2=\Phi_{\mathrm{int}}\Phi_{\mathrm{ext}}$, where $\Phi_{\mathrm{int}}=\prod_s\Phi_s$ with $\Phi_s=\prod_{i<j\in I_s}(x_i-x_j)^2$, and $\Phi_{\mathrm{ext}}$ is the product of the factors $(x_i-x_j)^2$ with $x_i\ne x_j$. The functional $\bigotimes_i\ell_{x_i}$ is the sum over $D\subset\{1,\dots,h\}$ of $\prod_{i\notin D}\alpha_{x_i}\prod_{i\in D}\beta_{x_i}$ times evaluation at $\mathbf x$ of $\partial_D=\prod_{i\in D}\partial/\partial x_i$. The polynomial $\Phi_s$ is homogeneous of degree $m_s(m_s-1)$ in the differences $x_i-x_j$ ($i,j\in I_s$); any derivative of it of order at most $m_s$ in these variables vanishes at the diagonal point unless its degree $m_s(m_s-1)-(\text{order})$ is zero. Hence $\mathfrak T(\mathbf x)=0$ if some $m_s\ge3$, and if $m_s=2$ only the terms with $I_s\subset D$ in which both derivatives fall on $\Phi_s$ survive, with $\partial_{x_i}\partial_{x_j}(x_i-x_j)^2=-2$. Writing $\mathsf S_1$ for the indices at singly occupied sites,
\[
\mathfrak T(\mathbf x)=\prod_{m_s=2}\bigl(-2\beta_s^2\bigr)\sum_{D\subset\mathsf S_1}\prod_{i\in\mathsf S_1\setminus D}\alpha_{x_i}\prod_{i\in D}\beta_{x_i}\,(\partial_D\Phi_{\mathrm{ext}})(\mathbf x).
\]
By Leibniz's rule each derivative $\partial/\partial x_i$ falls on one of the factors of $\Phi_{\mathrm{ext}}$ that contain $x_i$, of which there are fewer than $h$; so $\partial_D\Phi_{\mathrm{ext}}$ is a sum of at most $h^{|D|}$ terms, in each of which every factor $(x_i-x_j)^2$ is differentiated at most twice (once in $x_i$ and once in $x_j$). The derivatives of $(x_i-x_j)^2$ are $\pm2(x_i-x_j)$ and $-2$. Since $\Phi_{\mathrm{ext}}$ contains only factors with $x_i\ne x_j$, and the $x_i$ are integers, $|x_i-x_j|\ge1$ for each of them, so each differentiated factor is at most $2(x_i-x_j)^2$ in absolute value. So $|\partial_D\Phi_{\mathrm{ext}}(\mathbf x)|\le(2h)^{|D|}|\Phi_{\mathrm{ext}}(\mathbf x)|$. Finally $\alpha_x/\beta_x=(xW_n)'/(xW_n)=x^{-1}+3\sum_{a\le A}(x+a)^{-1}-\sum_{l\le K}(x+l)^{-1}$, so $|\alpha_x|\le(4+4\log K)|\beta_x|$. Summing over $D$ gives the bound with $C_0h\ge2(4+4\log K+2h)$.
\end{proof}

\begin{lemma}[the weight on the integers]\label{lem:weight-positive}
For $y\ge0$ let
\[
V(y)=3(y+a)\log(y+a)-2y\log y-(y+\kappa)\log(y+\kappa)+d-ty ,
\]
with $0\log0=0$ and $t=-\log|\lambda|$. There is $C_1$ such that for all $n$ and all integers $x\ge1$, with $y=x/h$,
\[
\log|\beta_x|\le-dh\log h+hV(y)+C_1\log(h+x).
\]
\end{lemma}

\begin{proof}
We have $\log|\beta_x|=-tx+\log x+3\log D_A(x)-\log D_K(x)$. Since $s\mapsto\log(x+s)$ is increasing, $\log D_m(x)\le\int_1^{m+1}\log(x+s)\dd s$ and $\log D_m(x)\ge\int_0^m\log(x+s)\dd s$. Therefore
\[
3\log D_A(x)-\log D_K(x)\le3\int_0^A\log(x+s)\dd s-\int_0^K\log(x+s)\dd s+3\log(x+A+1).
\]
Substituting $x=hy$ and $s=hs'$, $\int_0^{\mu h}\log(x+s)\dd s=\mu h\log h+h\bigl(F(y+\mu)-F(y)\bigr)$ with $F(z)=z\log z-z$. Collecting terms, $3\int_0^A-\int_0^K=-dh\log h+h\bigl(V(y)+ty\bigr)$.
\end{proof}

\subsection{The bound for positive arguments}\label{sec:positive-bound}

\begin{proposition}[comparison bound, positive arguments]\label{prop:positive-bound}
Let $\lambda\in\{2/11,2/9\}$, let $\sigma$ be a comparison measure supported in $[0,\infty)$, write $U=U_\sigma$, $E=E(\sigma)$ and $\mathsf W=V+2U$. Then for every $\theta\in\R$,
\[
\limsup_{n\to\infty}\frac{\log|\Qt_n(\xi)|}{h^2}\ \le\ \Cst-E+\theta+2\int_0^\infty\bigl(\mathsf W(y)-\theta\bigr)_+\dd y .
\]
\end{proposition}

\begin{proof}
By \eqref{eq:heine-sum} and \cref{lem:multiplicity}, $|Q_n(\xi)|\le\sum_{\mathbf x}|\mathfrak T(\mathbf x)|$, the sum running over tuples in which no value occurs more than twice. Fix such a tuple, put $y_i=x_i/h$, and note that $n_c\le2h$. Since $\log h>0$ and there are at most $h^2$ ordered pairs,
\[
\sum_{\substack{i\ne j\\x_i\ne x_j}}\log|x_i-x_j|\le h^2\log h+\sum_{\substack{i\ne j\\y_i\ne y_j}}\log|y_i-y_j|\le h^2\log h-h^2E+2h\sum_iU(y_i)+O(h\log h)
\]
by \cref{lem:energy-comparison}. With \cref{lem:multiplicity,lem:weight-positive},
\begin{equation}\label{eq:term-bound}
\log|\mathfrak T(\mathbf x)|\le(1-d)h^2\log h-h^2E+h\sum_i\mathsf W(y_i)+C_1\sum_i\log(h+x_i)+O(h\log h).
\end{equation}
As $y\to\infty$, $V(y)=-ty-d\log y+O(1/y)$ and $U(y)=\log y+O(1/y)$, so there are $Y$ and $c_1>0$ with $\mathsf W(y)-\theta\le-c_1(y-Y)$ for $y\ge Y$; hence $\mathsf W-\theta\le(\mathsf W-\theta)_+-c_1(y-Y)_+$ everywhere. Since each value occurs at most twice among the $x_i$,
\[
\sum_i\mathsf W(y_i)\le h\theta+2\sum_{x\ge1}\bigl(\mathsf W(x/h)-\theta\bigr)_+-\frac{c_1}h\sum_i(x_i-Yh)_+ .
\]
Insert this into \eqref{eq:term-bound}. The resulting bound for $|\mathfrak T(\mathbf x)|$ was derived for the tuples in which no value occurs more than twice; it is a positive quantity defined for every $\mathbf x\in\Z_{\ge1}^h$, so we may sum it over the admissible tuples and then enlarge the sum to all of $\Z_{\ge1}^h$, where it factorizes. We obtain
\begin{multline*}
\log|Q_n(\xi)|\le(1-d)h^2\log h+h^2(\theta-E)+2h\sum_{x\ge1}\bigl(\mathsf W(x/h)-\theta\bigr)_+\\
+h\log\Bigl(\sum_{x\ge1}(h+x)^{C_1}e^{-c_1(x-Yh)_+}\Bigr)+O(h\log h).
\end{multline*}
The last sum is $O(h^{C_1+1})$. The function $(\mathsf W-\theta)_+$ is continuous on $[0,\infty)$ with compact support, so $h^{-1}\sum_{x\ge1}(\mathsf W(x/h)-\theta)_+\to\int_0^\infty(\mathsf W-\theta)_+$. Adding \cref{lem:stirling} proves the proposition.
\end{proof}

\begin{remark}[scale for $\lambda=2/11$]\label{rem:211-scale}
For $\lambda=2/11$ the certificate is formulated in the variable $s=x/n=4y$ and divides by $n^2=h^2/16$. The substitution $y=s/4$ transforms \cref{prop:positive-bound} into
\[
\limsup_{n\to\infty}\frac{\log|\Qt_n(\xi)|}{n^2}\le24\log6-16\log4+12-E'+4\theta'+2\int_0^\infty\bigl(\mathsf W'(s)-\theta'\bigr)_+\dd s,
\]
where $\sigma'$ is the image of $4\sigma$ under $y\mapsto4y$ (a measure of mass $4$), $E'=\iint\log|s-s'|\dd\sigma'\dd\sigma'$, $\mathsf W'=V'+2U_{\sigma'}$ and $V'(s)=3(s+1)\log(s+1)-2s\log s-(s+6)\log(s+6)+3-ts$. Indeed $E'=16E+16\log4$, $\mathsf W'(4y)=4\mathsf W(y)+5\log4$, and with $\theta'=4\theta+5\log4$ the right side equals $16$ times that of \cref{prop:positive-bound}.
\end{remark}

\section{Archimedean decay II: a contour version for negative arguments}\label{sec:contour}

For $\lambda<0$ the terms of the series in \cref{lem:series} alternate in sign. Taking absolute values term by term, as in \cref{sec:positive-bound}, discards this cancellation. We keep the sign by writing the series as a contour integral before estimating, as is done in \cite[Sections~6.1 and~6.3]{Litsey2026Li2} with polygonal contours through $0$ (``bent frames''); the contour used here consists of two rays from the point $\frac12$. The result, \cref{prop:negative-bound}, is the analogue of \cref{prop:positive-bound} for $\lambda\in\{-2/9,-2/7\}$ and is used in \cref{sec:proof-arch}. The statements are candidates in the sense of \cref{sec:intro-status}.

Throughout this section $\lambda=-e^{-t}$ with $t=\log(b/2)$, and $c=e^{\mathrm i\varphi}$ with $0<\varphi<\pi/2$; we take $c=(4+3\mathrm i)/5$ for $\lambda=-2/9$ and $c=(21+20\mathrm i)/29$ for $\lambda=-2/7$. Let $\Gamma$ be the path that runs along the ray $\{\frac12+cy:y\ge0\}$ from infinity to $\frac12$ and then along the ray $\{\frac12+\bar cy:y\ge0\}$ from $\frac12$ to infinity, and put
\[
k(z)=\frac{\pi\,e^{-tz}}{\sin\pi z},\qquad T=t\cos\varphi+\pi\sin\varphi .
\]

\begin{lemma}\label{lem:kernel-bound}
There is a constant $C_0$ such that $|k(z)|+|k'(z)|\le C_0\,e^{-t\real z-\pi|\imag z|}$ for all $z\in\Gamma$.
\end{lemma}

\begin{proof}
We have $k'=-k\,(t+\pi\cot\pi z)$ and $|e^{-tz}|=e^{-t\real z}$. If $|\imag z|\ge1$, then $|\sin\pi z|\ge\sinh(\pi|\imag z|)\ge\frac12(1-e^{-2\pi})e^{\pi|\imag z|}$ and $|\cot\pi z|\le\coth\pi$. The part of $\Gamma$ with $|\imag z|\le1$ is compact and meets the real axis only at $\frac12$, so $\sin\pi z$ does not vanish on it and $k,k'$ are bounded there.
\end{proof}

\begin{lemma}[contour representation]\label{lem:contour}
For every $R\in\cR$,
\[
\Lambda_\xi(R)=-\frac1{2\pi\mathrm i}\int_\Gamma k'(z)\,zR(z)\dd z ,
\]
the integral being absolutely convergent.
\end{lemma}

\begin{proof}
Put $g=(zR)'$, a rational function whose poles lie at negative integers. For $N\ge1$ let $\Delta_N$ be the triangle with vertices $\frac12$ and $N+\frac12\pm\mathrm iN\tan\varphi$. Its positively oriented boundary consists of the segment of the lower ray from $\frac12$ outwards, the vertical segment $\real z=N+\frac12$ upwards, and the segment of the upper ray back to $\frac12$. Inside $\Delta_N$ the only poles of $kg$ are the integers $1,\dots,N$, where $k$ has residue $e^{-tm}\cos(\pi m)^{-1}=\lambda^m$. By the residue theorem
\[
\sum_{m=1}^N\lambda^mg(m)=\frac1{2\pi\mathrm i}\oint_{\partial\Delta_N}k(z)g(z)\dd z .
\]
On the vertical side, $|\sin\pi z|=\cosh(\pi\imag z)\ge1$, so $|k(z)|\le\pi e^{-t(N+1/2)}$, while $|g(z)|$ grows at most polynomially in $N$ and the side has length $O(N)$; its contribution tends to $0$. By \cref{lem:kernel-bound} the integrals along the rays converge absolutely as $N\to\infty$, and with \cref{lem:series} we obtain $\Lambda_\xi(R)=(2\pi\mathrm i)^{-1}\int_\Gamma k\,(zR)'\dd z$. Integrating by parts along the single path $\Gamma$, the boundary terms at its two ends vanish because $k(z)zR(z)\to0$ exponentially, which gives the stated formula.
\end{proof}

\begin{lemma}[the weight on the contour]\label{lem:weight-negative}
For $y\ge0$ let
\begin{align*}
V_-(y)&=3\int_0^a\log|cy+s|\dd s-\int_0^\kappa\log|cy+s|\dd s-Ty\\
&=\real\bigl[3(cy+a)\log(cy+a)-2cy\log(cy)-(cy+\kappa)\log(cy+\kappa)\bigr]+d-Ty,
\end{align*}
with the principal branch and $0\log0=0$. There is $C_1$ such that for all $n$, all $y\ge0$ and $z=\frac12+hv$ with $v\in\{cy,\bar cy\}$,
\[
\log\bigl|zW_n(z)k'(z)\bigr|\le-dh\log h+hV_-(y)+C_1\log(h+hy).
\]
\end{lemma}

\begin{proof}
By \cref{lem:kernel-bound}, $\log|k'(z)|\le\log C_0-t\real z-\pi|\imag z|\le\log C_0-hTy$, and $\log|z|\le\log(1+hy)$. Since $\real z>0$, the function $s\mapsto|z+s|$ is increasing on $[0,\infty)$, so as in \cref{lem:weight-positive}
\[
3\log|D_A(z)|-\log|D_K(z)|\le3\int_1^{A+1}\log|z+s|\dd s-\int_0^K\log|z+s|\dd s .
\]
For $s>0$ the number $w=hv+s$ has $\real w\ge s$, hence $|w|\le|w+\frac12|\le|w|(1+\frac1{2s})$, that is $\log|hv+s|\le\log|z+s|\le\log|hv+s|+\log(1+\frac1{2s})$. Using the upper bound in the first integral and the lower bound in the second, and $\int_1^{A+1}\log(1+\frac1{2s})\dd s=O(\log h)$, we may replace $z$ by $hv$ at the cost $O(\log h)$. Replacing $\int_1^{A+1}$ by $\int_0^A$ costs at most $3\log(h|v|+A+1)+3$. Finally $\int_0^{\mu h}\log|hv+s|\dd s=\mu h\log h+h\int_0^\mu\log|v+s'|\dd s'$ and $|\bar cy+s'|=|cy+s'|$. The closed form follows from $\int_0^\mu\log|v+s|\dd s=\real\bigl[(v+\mu)\log(v+\mu)-v\log v\bigr]-\mu$.
\end{proof}

\begin{proposition}[comparison bound, negative arguments]\label{prop:negative-bound}
Let $\lambda\in\{-2/9,-2/7\}$ and let $\sigma$ be a comparison measure supported on the two rays $\{cy:y\ge0\}$ and $\{\bar cy:y\ge0\}$ and invariant under complex conjugation. Write $U_-(y)=U_\sigma(cy)=U_\sigma(\bar cy)$ and $E=E(\sigma)$. Then
\[
\limsup_{n\to\infty}\frac{\log|\Qt_n(\xi)|}{h^2}\ \le\ \Cst-E+\sup_{y\ge0}\bigl(V_-(y)+2U_-(y)\bigr).
\]
\end{proposition}

\begin{proof}
By \cref{lem:contour}, $\Lambda_\xi(W_nx^s)=\int_\Gamma z^s\dd\nu(z)$ for the complex measure $\dd\nu=-(2\pi\mathrm i)^{-1}k'(z)zW_n(z)\dd z$ on $\Gamma$, all of whose moments converge absolutely. \Cref{lem:heine} and Fubini's theorem give
\[
Q_n(\xi)=\frac1{h!}\int_{\Gamma^h}\prod_{i<j}(z_j-z_i)^2\prod_i\dd\nu(z_i).
\]
Parametrize $z_i=\frac12+hv_i$ with $v_i\in\{cy_i,\bar cy_i\}$, $y_i\ge0$, so that $|\dd z_i|=h\dd y_i$ and $\prod_{i\ne j}|z_i-z_j|=h^{h(h-1)}\prod_{i\ne j}|v_i-v_j|$. Outside a null set the $v_i$ are distinct, and \cref{lem:energy-comparison} gives $\sum_{i\ne j}\log|v_i-v_j|\le-h^2E+2h\sum_iU_-(y_i)+O(h\log h)$; here $n_c=h$ and $m_\sigma=2$, and $U_\sigma(v_i)=U_-(y_i)$ whichever of the two rays carries $v_i$, because $\sigma$ is invariant under complex conjugation. This is why a single supremum over $y\ge0$ controls both rays. Let $\theta$ be the supremum in the statement; it is finite because $V_-(y)=-Ty+O(\log(2+y))$ and $U_-(y)=O(\log(2+y))$, and for the same reason there are $Y$ and $c_1>0$ with $V_-+2U_-\le\theta-c_1(y-Y)_+$. With \cref{lem:weight-negative} the integrand is bounded by
\[
\exp\Bigl((1-d)h^2\log h-h^2E+h^2\theta+O(h\log h)\Bigr)\prod_i(h+hy_i)^{C_1}e^{-hc_1(y_i-Y)_+},
\]
and each of the $h$ integrations over $\Gamma$ contributes a factor $O(h^{C_1+2})$. Hence $\log|Q_n(\xi)|\le(1-d)h^2\log h+h^2(\theta-E)+O(h\log h)$, and \cref{lem:stirling} completes the proof.
\end{proof}

\begin{remark}
The choice of the rays is a compromise. A steeper ray increases the decay rate $T$ of the kernel but moves the points away from the real axis, where the weight $W_n$ and the energy are less favourable. The directions $c$ used here have rational real and imaginary parts, which keeps the certificates in exact rational arithmetic; any direction with a valid certificate would do. We have not compared the resulting bounds with those of the term-by-term method of \cref{sec:positive-bound} applied to $|\lambda|$.
\end{remark}

\section{Proof of the theorems}\label{sec:proof}

This section completes the proofs of \cref{thm:arch} and \cref{thm:main}, discusses how robust the numerical margins are, and lists the open points. \Cref{thm:arith,thm:nonvanishing} were proved in \cref{sec:profile-values,sec:reduction}.

\subsection{The archimedean estimate}\label{sec:proof-arch}
\begin{proof}[Proof of \cref{thm:arch}]
For each $\lambda$, \cref{app:certificates} specifies a piecewise constant comparison measure $\sigma_\lambda$ and, for the positive arguments, a number $\theta_\lambda$. For $\lambda=2/9$, \cref{prop:positive-bound} applies directly; for $\lambda=2/11$ it applies in the form of \cref{rem:211-scale}; for $\lambda=-2/9,-2/7$, \cref{prop:negative-bound} applies. The certificates bound the right-hand sides from above by the values of \cref{tab:analytic}, which are below the thresholds $\ran(\lambda)$ of \cref{tab:rates}.
\end{proof}

\begin{table}[ht]
\centering
\caption{The computed upper bounds for the comparison-measure functionals, normalized by $h^2$. For $\lambda>0$, ``penalty'' is the computed upper bound for $\int_0^\infty(\mathsf W-\theta)_+$; for $\lambda<0$, $\theta$ is the computed upper bound for $\sup_{y\ge0}(V_-+2U_-)$. The last column is $\Cst-E+\theta+2\cdot\text{penalty}$, respectively $\Cst-E+\theta$, rounded up; the other entries are rounded to the nearest digit. For $\lambda=2/11$ the certificate works in the scale $n^2=h^2/16$, with $E'=4.4615788$, $\theta'=-431/50$, penalty bound $0.6028348$ and bound $-4.9143918$; the row shows these values converted by \cref{rem:211-scale}, $E=E'/16-\log4$, $\theta=(\theta'-5\log4)/4$ and penalty${}/16$. There $\int_0^\infty(\mathsf W'-\theta')_+\approx0.35243$, and the bound $0.6028348$ includes the bounded-variation error term of the midpoint rule. For $\lambda=2/9$ the certificate program does not print $E(\sigma)$ and the penalty bound separately; the table shows the value of $E(\sigma)$ from the recomputation described in the section on reproducibility and the penalty bound implied by the computed total.}
\label{tab:analytic}
\begin{tabular}{@{}lccccc@{}}
\toprule
$\lambda$ & $\Cst$ & $E(\sigma)$ & $\theta$ & penalty & bound\\
\midrule
$2/11$ & $2.3979184$ & $-1.1074457$ & $-3.8878680$ & $0.0376772$ & $-0.3071494$\\
$2/9$ & $0.5075059$ & $-1.3038987$ & $-167/50$ & $0.0924819$ & $-1.3436316$\\
$-2/9$ & $2.5190484$ & $-1.3538434$ & $-4.3683007$ & --- & $-0.4954089$\\
$-2/7$ & $2.2176686$ & $-1.3534233$ & $-4.0971701$ & --- & $-0.5260781$\\
\bottomrule
\end{tabular}
\end{table}

\subsection{Proof of the main theorem}\label{sec:proof-main}
\begin{proof}[Proof of \cref{thm:main}]
Fix $\lambda$ and let $p_k$ be the $k$th prime. Apply \cref{lem:criterion} with $x=\xi$, $F_k=\Qt_{p_k}$, $h_k=Hp_k$, $\rho=\rar(\lambda)$ and $\rho'=\ran(\lambda)$; then $\deg F_k\le h_k$ by \cref{sec:determinants}, and $\rho>\rho'$ by \cref{tab:rates}. The nonvanishing hypothesis is \cref{thm:nonvanishing}. The content hypothesis follows from \cref{thm:arith}, since $F_k\ne0$ for large $k$, and the value hypothesis from \cref{thm:arch}. Hence $\xi$ is irrational.
\end{proof}

\subsection{Margins and their robustness}\label{sec:margins}
The margins of \cref{tab:rates} between the computed arithmetic and analytic bounds are $0.036084$ (that is, $0.577$ in the $n^2$-scale), $0.019104$, $0.112321$ and $0.0047556$; the margins between the thresholds $\rar$ and $\ran$, which are the numbers that enter \cref{lem:criterion}, are $0.035125$, $0.010$, $0.090$ and $0.004755$. They are smaller than what the same comparison measures and profiles actually give, because the certificates use simple error bounds. The AI-internal recomputation described in the section on reproducibility encloses the same quantities rigorously on finer meshes, and evaluates the profile integrals exactly.

For the negative arguments the certificates bound a supremum on a mesh of width $1/400$ and add a uniform interpolation error, $0.026348$ for $-2/9$ and $0.02389$ for $-2/7$. With the same comparison measures the recomputed bounds for the right side of \cref{prop:negative-bound} are $-0.52162$ and $-0.54985$, instead of $-0.49541$ and $-0.52608$; in both cases the supremum of $V_-+2U_-$ is attained at the vertex $y=0$. The midpoint lower bounds for the profile integrals are also conservative there: the exact values of $I_\delta$ are $0.168559$ and $0.078109$, instead of $0.167309$ and $0.076945$. The rigorous margins for $-2/9$ and $-2/7$ are therefore at least $0.13978$ and $0.02969$; the thin certified margin for $-2/7$ reflects the error bounds, not the comparison measure. For $\lambda=2/11$ the bounded-variation estimate is similarly generous: the penalty integral is about $0.352$ in the $n^2$-scale, against the certified $0.603$, and the right side of \cref{rem:211-scale} is at most $-5.4152$ rather than $-4.9144$; the margin becomes $1.084$ in the $n^2$-scale, or $0.0678$ per $h^2$.

For $\lambda=2/9$ the situation is different, and this case is the fragile one. The recomputed bound for the right side of \cref{prop:positive-bound} is $-1.34391$, within $3\cdot10^{-4}$ of the computed $-1.34363$, and the profile integral is evaluated exactly by the midpoint rule because $\mathsf B$ is affine on every subinterval; the recomputed margin is $0.019379$, against rates of size $1.3$. The analytic bound cannot be improved much with this comparison measure. The exact data do not support it: at the largest computed index, $n=11$, the value $h^{-2}\log|\Qt_n(\xi)|=-1.2886$ lies above the upper bound $-1.3436$ for the upper limit, and also above $\rar=-1.33$ (\cref{tab:numerics}); the sequence is still decreasing, and extrapolations of it give limits between $-1.349$ and $-1.341$, an interval that contains the bound. For the other three values the finite data lie below $\ran$. So for $2/9$ the theorem requires the upper bound of \cref{sec:energy} to be asymptotically sharp to within about $1.4\%$ and the lower bounds of \cref{sec:exceptional,sec:ordinary,sec:profile} to lose nothing beyond what is computed: any term of size $0.02\,h^2$ that is unaccounted for in these sections is fatal for this case, and nothing but the written proofs protects it. We have neither extended the exact computation far enough to see the sequence cross the bound (an attempt at $n=12,\dots,15$ was stopped for lack of time, without result), nor found a triple $(H,B,C)$ with a larger margin, nor solved the constrained equilibrium problem by a second method.

\subsection{Open points}\label{sec:open}
We list the points of the argument that we regard as open, or as deserving particular scrutiny by a reader.
\begin{enumerate}[(1)]
\item \emph{External check of the certificates.} The profile integrals of \cref{prop:profile-certificate} and the comparison-measure bounds of \cref{tab:analytic} have been reproduced by the original programs and recomputed, with confirmed results, by a separately written implementation; the eight local determinants have been recomputed by a third program. All these implementations were written with AI systems, and the recomputations are AI-internal checks. A check by an independent human or external party is outstanding.
\item \emph{Error bounds inside the certificates.} The certificates rely on explicit error estimates: the weak-duality bound and the truncation of jet levels (\cref{sec:cert-profile}), the bounded-variation midpoint rule for $\lambda=2/11$, the $L^1$ interpolation bound for $\lambda=2/9$, the uniform interpolation bounds and the tail estimates for the negative arguments (\cref{sec:cert-measures}). Each is proved in \cref{app:certificates}, but these proofs have not been reviewed independently, and the small margin for $2/9$ makes them consequential.
\item \emph{The distribution argument.} \Cref{prop:distribution,lem:local-integrality} replace $p$-adic integration by finite normalized sums and an affine change of variable. They are used in the jet and pole bounds and in \cref{thm:nonvanishing-general}. The exact computations of \cref{app:numerics} test their consequences, not the lemmas themselves.
\item \emph{Constants.} The $O(h\log h)$ terms of \cref{sec:exceptional,sec:energy,sec:contour} and the rate of convergence in \cref{lem:pnt} are not made explicit. This is harmless for \cref{thm:main}, but a formal proof, or an effective irrationality measure, would need them.
\item \emph{Review and formalization.} No independent human review of the general arguments has taken place, and \cref{thm:arith,thm:arch,thm:nonvanishing} are not formalized.
\item \emph{Literature.} The method is that of Tang and Litsey with different parameters (\cref{sec:intro-sources}). We have compared our arguments with \cite{Litsey2026Li2} at the level of statements only, and we make no claim of novelty for the techniques. The sources of 2026 are unrefereed, and the threshold quoted for the refereed literature is taken second-hand (\cref{sec:intro-context}).
\item \emph{The case $2/9$.} See \cref{sec:margins}: a longer exact computation, a second solution of the equilibrium problem, or a family with a larger margin would each strengthen this case.
\end{enumerate}


\appendix
\section{The certificates}\label{app:certificates}

This appendix describes the finite computations used in \cref{prop:profile-certificate,prop:local-dets} and in the proof of \cref{thm:arch}, and proves the error estimates on which they rely; these computations are the certified steps of \cref{tab:status}. All assertions are checked in exact rational arithmetic; decimal values are for orientation only.

\subsection{Rational enclosures of logarithms and arctangents}\label{sec:cert-intervals}
For a rational $z\in[1,2]$ put $r=(z-1)/(z+1)\in[0,\frac13]$. Since $\log z=2\operatorname{artanh}r$,
\[
2\sum_{k=0}^{19}\frac{r^{2k+1}}{2k+1}\ \le\ \log z\ \le\ 2\sum_{k=0}^{19}\frac{r^{2k+1}}{2k+1}+\frac{2r^{41}}{41(1-r^2)},
\]
the remainder being bounded by a geometric series. A positive rational number is reduced to $[1,2]$ by a power of $2$, using the same enclosure for $\log2$, and the endpoints are rounded outwards to multiples of $10^{-15}$. For $|u|\le\frac12$ the alternating series of $\arctan u$ through degree $39$ has error at most $|u|^{41}/41$; for $\frac12<u\le1$ one uses $\arctan u=\frac\pi4+\arctan\frac{u-1}{u+1}$, for $u>1$ one uses $\arctan u=\frac\pi2-\arctan\frac1u$, and oddness for $u<0$. The number $\pi$ is enclosed by Machin's formula $\pi=16\arctan\frac15-4\arctan\frac1{239}$. For a nonzero rational point $w$ of the closed upper half-plane, $\real\log w=\frac12\log|w|^2$ and $\imag\log w=\arg w$ is obtained from the arctangent by the quadrant rule. Interval arithmetic with rational endpoints then encloses every quantity below.

\subsection{The profile integrals}\label{sec:cert-profile}
Let $\mathfrak P$ be the partition of $[\delta,\kappa]$ by $\delta$, $\kappa$ and the points $s/j\in(\delta,\kappa)$ with $s\in\{a,\kappa,1,\kappa-a\}$ and $j\ge1$; for $\lambda=2/11$, in the scale $y=4x$, by the points $s/j$ with $s\in\{1,6,4,5\}$ in $(\frac1{40},6)$. The numbers of subintervals are $479$, $999$, $1749$ and $1799$ for $\lambda=2/11$, $2/9$, $-2/9$ and $-2/7$. By \cref{lem:profile-regular}, on each subinterval $\mathcal I$ the weights of all slots are constant and their capacities are affine in $x$.

\begin{lemma}\label{lem:profile-convex}
On each subinterval $\mathcal I$ of $\mathfrak P$ the function $\mathsf B$ is convex. Consequently $\int_{\mathcal I}\mathsf B\ge|\mathcal I|\,\mathsf B(x_{\mathcal I})$, where $x_{\mathcal I}$ is the midpoint of $\mathcal I$.
\end{lemma}

\begin{proof}
The minimum of a linear program $\min\{w\cdot y:\ 0\le y\le c,\ \sum y=1\}$, or with $\sum y\le1$, is a convex function of the capacity vector $c$, since a convex combination of feasible points for $c$ and $c'$ is feasible for the convex combination of $c$ and $c'$. Composed with the affine map $x\mapsto c(x)$, this shows that $J$ and $R$ are convex on $\mathcal I$; so is their maximum, and $\mathsf N$ is affine on $\mathcal I$. The last assertion is Jensen's inequality.
\end{proof}

The certificate evaluates $\mathsf B(x_{\mathcal I})$ exactly for every subinterval, filling the cheapest slots first, with jet levels $j<q_1+3$ (which suffice, as noted after \cref{def:profile}), and sums $|\mathcal I|\mathsf B(x_{\mathcal I})$. As a second justification, a separate program computes on each subinterval an explicit affine minorant of $\max(J,R)$ from the weak-duality bound of \cref{lem:weak-duality}, with a price chosen at the midpoint and held fixed; it reproduces the midpoint values exactly on all $5026$ subintervals. The resulting rational lower bounds for $I_\delta$ are those of \cref{tab:profile}. Finally the exceptional contribution $\eps_{\mathrm{exc}}$ is enclosed with the logarithms of \cref{sec:cert-intervals}, and the inequality $I_\delta+\eps_{\mathrm{exc}}>\rar$ is checked.

\subsection{The local determinants}\label{sec:cert-local}
For $s<2H-1$ the function $u^sw_b$ is written as $\rho+\sum_k\varrho_k(u+k)^{-1}$ by division by the monic denominator and $\varrho_k=\mathrm{num}(-k)/\prod_{k'\ne k}(k'-k)$; then $\cV(u^sw_b)=\sum_{m\ge1}m[u^m]\rho\,M_{m-1}(\lambda)+\sum_k\varrho_k\lambda^{-k}\tau_k(\lambda)$, with $M_k$ from \eqref{eq:moments}. The determinants of the $H\times H$ Hankel matrices are computed exactly, and each is certified a second time by an explicit factorization $\mathsf H_b=\mathsf L\mathsf U$ with $\mathsf L$ unit lower triangular and $\mathsf U$ upper triangular, checked entrywise. A separate program, which computes the moments by the recursion $(\lambda-1)M_k=\sum_{j<k}\binom kj(-1)^{k-j}M_j$ instead of Stirling numbers, reproduces all eight values of \cref{tab:local-dets}.

\subsection{The comparison measures}\label{sec:cert-measures}
All comparison measures are piecewise constant in a parameter $y\ge0$ on a uniform grid. Put
\[
f(s)=s\log|s|-s,\qquad g(s)=\tfrac12s^2\log|s|-\tfrac34s^2\qquad(f(0)=g(0)=0),
\]
so that $g'=f$ and $f'=\log|\cdot|$. For the unit density on $[\alpha,\beta]\subset\R$, the potential at $y\in\R$ is $f(y-\alpha)-f(y-\beta)$, and the interaction with the unit density on $[\gamma,\gamma']$ is
\[
\int_\alpha^\beta\!\!\int_\gamma^{\gamma'}\log|s-s'|\dd s'\dd s=g(\beta-\gamma)+g(\alpha-\gamma')-g(\alpha-\gamma)-g(\beta-\gamma').
\]
For the negative arguments the measure lives on the two rays $\{cy\}$ and $\{\bar cy\}$ and is symmetric. With $\eta=c^2$ one has $|cy-\bar cs|=|\eta y-s|$. Put $F(z)=z(\log z-1)$ and $G(z)=z^2(\frac12\log z-\frac34)$ on the closed upper half-plane (principal branch, $F(0)=G(0)=0$), so that $\real F=f$ and $\real G=g$ on the real axis and $G''=\log$. Then the cross-ray potential at $cy$ of the unit density on $[\alpha,\beta]$ is $\real\bigl(F(\eta y-\alpha)-F(\eta y-\beta)\bigr)$, and the cross-ray interaction of the unit densities on $[\alpha,\beta]$ and $[\gamma,\gamma']$ is $\real\bigl[\eta^{-1}\bigl(G(\eta\beta-\gamma)+G(\eta\alpha-\gamma')-G(\eta\alpha-\gamma)-G(\eta\beta-\gamma')\bigr)\bigr]$. The energy is the sum of the same-ray and cross-ray interactions, each counted twice.

The data are as follows; the integer vectors $(m_j)$ are listed in \cref{tab:mass-vectors}.
\begin{itemize}
\item $\lambda=2/11$ (scale $s=4y$, mass $4$): density $m_j/500$ on $[\frac j2,\frac{j+1}2]$, $0\le j\le16$, $\sum m_j=4000$; $\theta'=-431/50$.
\item $\lambda=2/9$: density $m_j/2500$ on $[\frac j{40},\frac{j+1}{40}]$, $0\le j\le76$, $\sum m_j=10^5$; $\theta=-167/50$.
\item $\lambda=-2/9$, $c=(4+3\mathrm i)/5$: on each ray, density $m_j/5000$ on $[\frac j{40},\frac{j+1}{40}]$, $0\le j\le38$, $\sum m_j=10^5$.
\item $\lambda=-2/7$, $c=(21+20\mathrm i)/29$: on each ray, density $m_j/8000$ on $[\frac j{25},\frac{j+1}{25}]$, $0\le j\le23$, $\sum m_j=10^5$.
\end{itemize}

\begin{table}[ht]
\centering
\caption{The mass vectors $(m_j)$ of the comparison measures.}
\label{tab:mass-vectors}
\footnotesize
\begin{tabular}{@{}lp{0.8\textwidth}@{}}
\toprule
$2/11$ & 1000, 557, 402, 325, 272, 233, 203, 178, 157, 139, 122, 107, 92, 78, 63, 48, 24\\
$2/9$ & 5000 (seven times), 3819, 3235, 2875, 2603, 2389, 2213, 2065, 1937, 1826, 1728, 1641, 1562, 1490, 1425, 1365, 1310, 1259, 1211, 1167, 1125, 1086, 1049, 1014, 981, 950, 920, 891, 864, 837, 812, 788, 764, 742, 720, 699, 678, 659, 639, 620, 602, 584, 567, 549, 533, 516, 500, 484, 468, 452, 437, 421, 406, 391, 376, 361, 346, 330, 315, 300, 284, 268, 252, 235, 217, 199, 179, 158, 134, 107, 71\\
$-2/9$ & 9424, 5610, 5336, 4824, 4487, 4204, 3966, 3757, 3571, 3403, 3248, 3106, 2972, 2847, 2729, 2616, 2508, 2405, 2305, 2208, 2114, 2023, 1933, 1844, 1757, 1671, 1585, 1499, 1413, 1326, 1238, 1148, 1055, 958, 856, 744, 618, 477, 215\\
$-2/7$ & 13112, 8416, 7787, 6951, 6348, 5833, 5390, 4997, 4642, 4318, 4016, 3734, 3466, 3210, 2962, 2721, 2482, 2244, 2002, 1751, 1486, 1185, 869, 78\\
\bottomrule
\end{tabular}
\end{table}

\emph{The penalty integral for $\lambda=2/11$.} Let $\mathsf W'=V'+2U_{\sigma'}$ as in \cref{rem:211-scale}. For $s\ge12$ the support $[0,\frac{17}2]$ of $\sigma'$ gives $U_{\sigma'}(s)\le4\log s$, and $\frac{\mathrm d}{\mathrm ds}(V'+8\log s)=3\log(1+\frac1s)-\log(1+\frac6s)+\frac8s-t\le\frac{11}s-t<-\frac12$ because $t=\log\frac{11}2>\frac{17}{12}$; the certificate checks $V'(12)+8\log12<\theta'$, so $(\mathsf W'-\theta')_+=0$ on $[12,\infty)$. On $[0,12]$ the total variation of $\mathsf W'$ is at most that of $V'$, bounded by $\int_0^{12}\bigl(3\log(1+\frac1s)+\log(1+\frac6s)+t\bigr)\dd s$, plus twice the sum over the cells of the density times the variation of the cell potential; the potential of a cell decreases up to the midpoint of the cell and increases afterwards, so its variation on $[0,12]$ is the sum of its endpoint values minus twice its midpoint value. The certificate finds a total variation below $101$. For a function of bounded variation the midpoint rule with step $s_0$ has error at most $\frac{s_0}2$ times the variation, and $(\cdot-\theta')_+$ does not increase variation. With $s_0=\frac1{200}$ and upper enclosures of the midpoint values, the midpoint sum is $0.35241$ and the integral itself is about $0.35243$; adding the error term, $\frac{s_0}2$ times the computed total-variation bound, gives the certified bound $\int_0^\infty(\mathsf W'-\theta')_+<0.602835$ used in the certificate.

\emph{The penalty integral for $\lambda=2/9$.} For $y\ge4$, $U\le\log y$ and $V+2\log y$ has derivative $\log\frac y{y+\kappa}+\frac2y-t<0$; the certificate checks $V(4)+2\log4<\theta$. On $[0,4]$ let $\mathsf L$ be the piecewise linear interpolant of $\mathsf W$ on the mesh of width $s_0=\frac1{400}$; all breakpoints of $\sigma$ and the shift $-\kappa$ lie on the mesh. Since $(\cdot)_+$ is $1$-Lipschitz and the positive part of an affine function is convex, $\int_0^4(\mathsf W-\theta)_+\le\int_0^4(\mathsf L-\theta)_++\|\mathsf W-\mathsf L\|_{L^1}$, and the first integral is at most the trapezoidal sum of the positive parts of the (upper enclosures of the) nodal values. The function $\mathsf W$ is a combination of translates $f(\cdot-\alpha)$ with mesh-aligned $\alpha$, plus a linear function, with absolute coefficient sum $10$ (two from $V$, and $2\times4$ from $2U$, the density being nonincreasing with $\rho_0=2$). For one translate, the $L^1$ interpolation error on a cell adjacent to $\alpha$ is $s_0^2/4$ (by direct integration of $x\log x-x$), and on a cell at distance at least $js_0$ it is at most $s_0^3\max|f''|/12\le s_0^2/(12j)$. Summing over cells gives at most $s_0^2(\frac12+\frac16\sum_{j\le N}\frac1j)$ with $N=2050$, hence $\|\mathsf W-\mathsf L\|_{L^1}<1.17\cdot10^{-4}$.

\emph{The supremum for the negative arguments.} On $[0,Y]$, with $Y=2$ for $\lambda=-2/9$ and $Y=3$ for $\lambda=-2/7$, the function $V_-+2U_-$ is a linear function plus a combination of the functions $f(y-\alpha)$, $\real F(\eta y-\alpha)$ and $\real\bigl((cy+\beta)\log(cy+\beta)\bigr)$ with mesh-aligned $\alpha,\beta\ge0$ (mesh $s_0=\frac1{400}$), with absolute coefficient sum $6+4\mathsf J$, where $\mathsf J$ is the total variation of the one-ray density including its jumps at the ends of the support. Each of these functions differs from its linear interpolant on the mesh by at most $s_0/2$ uniformly: for $f(y-\alpha)$ the error is at most $s_0/e$ on the two cells adjacent to $\alpha$ and $s_0/8$ elsewhere; for $\real F(\eta y-\alpha)$ the second derivative is bounded by $|\eta y-\alpha|^{-1}$, and $|\eta y-\alpha|\ge s_0\sin(\arg\eta)$ except on the first cell when $\alpha=0$, where the function is $\real(\eta)\,y\log y$ plus a linear function; the factors $(cy+\beta)\log(cy+\beta)$ are treated in the same way. Hence $\sup_{[0,Y]}(V_-+2U_-)$ is at most the largest nodal upper enclosure plus $(6+4\mathsf J)s_0/2$, which is $0.026348$ for $-2/9$ and $0.02389$ for $-2/7$. For $y\ge Y$ one has $|cy+s|\ge y$ and $|cy+s|\le y+s$, so $V_-(y)+2U_-(y)\le3a\log(y+a)-\kappa\log y-Ty+2\log(y+b_*)$, where $b_*$ is the end of the support; this majorant decreases for $y\ge Y$ because $(3a+2)/Y<T$, and the certificate checks that its value at $Y$ is below the bound obtained on $[0,Y]$.

\section{Technical lemmas}\label{app:technical}

This appendix proves five general lemmas used in the body: a coefficient bound for integer-valued polynomials (used in \cref{lem:integral-valued}), a valuation bound for determinants of sums of rank-one matrices (used in \cref{sec:exceptional}), the divisibility of alternating polynomials by the Vandermonde product (used in \cref{prop:pole}), and two facts on linear programs with capacities (used in \cref{sec:profile,sec:cert-profile}).

\begin{lemma}[coefficients of integer-valued polynomials]\label{lem:integer-valued-coefficients}
Let $p$ be a prime, $N\ge1$, $L=\lfloor\log N/\log p\rfloor$, and let $g\in\Q[y]$ have degree at most $N$ and take values in $\Zp$ at all integers. Then for every $t\in\Z$ and $s\ge0$ the coefficient of $y^s$ in $g(t+y)$ has valuation at least $-sL$.
\end{lemma}

\begin{proof}
Newton's formula gives $g(t+y)=\sum_{k\le N}\Delta^kg(t)\binom yk$ with $\Delta^kg(t)\in\Zp$. For $k\ge1$,
\[
\binom yk=\frac{(-1)^{k-1}y}{k}\prod_{m=1}^{k-1}\Bigl(1-\frac ym\Bigr),
\]
so the coefficient of $y^s$ in $\binom yk$ is a sum of products of $s$ factors $\pm1/m$ with $m\le k\le N$, each of valuation at least $-L$.
\end{proof}

\begin{lemma}[slot lemma]\label{lem:slots}
Let $p$ be a prime and let $\mathsf M=\sum_{\varsigma\in\mathsf S}\mathbf u_\varsigma\mathbf w_\varsigma^{\mathsf T}$ be a finite sum of rank-one $h\times h$ matrices with $\mathbf u_\varsigma,\mathbf w_\varsigma\in\Q[X]^h$. Let $\omega_\varsigma\le\min_i\vg p(u_{\varsigma,i})+\min_j\vg p(w_{\varsigma,j})$. Then $\vg p(\det\mathsf M)\ge\min\sum_{\varsigma\in\mathsf S'}\omega_\varsigma$, the minimum over subsets $\mathsf S'\subset\mathsf S$ with $h$ elements.
\end{lemma}

\begin{proof}
$\mathsf M=\mathsf U\mathsf W^{\mathsf T}$ with $\mathsf U=[\mathbf u_\varsigma]$ and $\mathsf W=[\mathbf w_\varsigma]$, and by the Cauchy--Binet formula $\det\mathsf M=\sum_{\mathsf S'}\det\mathsf U_{\bullet\mathsf S'}\det\mathsf W_{\bullet\mathsf S'}$. Each term of the permutation expansion of $\det\mathsf U_{\bullet\mathsf S'}$ contains exactly one entry of each column, so $\vg p(\det\mathsf U_{\bullet\mathsf S'})\ge\sum_{\varsigma\in\mathsf S'}\min_i\vg p(u_{\varsigma,i})$, and similarly for $\mathsf W$.
\end{proof}

\begin{lemma}[alternating polynomials]\label{lem:alternating}
Let $p$ be an odd prime and $\Phi\in\Zp[z_1,\dots,z_k]$ alternating, that is, changing sign under each transposition of two variables. Then $\Phi$ is divisible by $\prod_{i<j}(z_i-z_j)$ in $\Zp[z_1,\dots,z_k]$.
\end{lemma}

\begin{proof}
Fix $i<j$. Division by the monic polynomial $z_i-z_j$ in the variable $z_i$ gives $\Phi=(z_i-z_j)\Psi+\Phi|_{z_i=z_j}$. The transposition of $z_i$ and $z_j$ fixes $\Phi|_{z_i=z_j}$ and changes the sign of $\Phi$, so $2\Phi|_{z_i=z_j}=0$, and $\Phi|_{z_i=z_j}=0$ because $2$ is a unit. Thus each $z_i-z_j$ divides $\Phi$. These are pairwise non-associate prime elements of the factorial ring $\Zp[z_1,\dots,z_k]$, so their product divides $\Phi$.
\end{proof}

\begin{lemma}[linear programs with capacities]\label{lem:lp-lipschitz}
Let $w_1,\dots,w_r$ be real numbers with $|w_k|\le W$, and for capacities $c\in[0,\infty)^r$ put
\begin{align*}
F(c)&=\min\Bigl\{\sum_kw_ky_k:\ 0\le y_k\le c_k,\ \sum_ky_k=m\Bigr\},\\
F^-(c)&=\min\Bigl\{\sum_kw_ky_k:\ 0\le y_k\le c_k,\ \sum_ky_k\le m\Bigr\}.
\end{align*}
If $\sum c_k\ge m$ and $\sum c'_k\ge m$, then $|F(c)-F(c')|\le2W\|c-c'\|_1$; and always $|F^-(c)-F^-(c')|\le W\|c-c'\|_1$. Both minima are attained by filling the slots in order of increasing weight (for $F^-$, only slots of negative weight).
\end{lemma}

\begin{proof}
Let $y$ be optimal for $c$ and put $y'_k=\min(y_k,c'_k)$. Then $\sum_k(y_k-y'_k)\le\|c-c'\|_1$. For $F^-$, $y'$ is feasible for $c'$ and costs at most $W\|c-c'\|_1$ more. For $F$, refill the deficit using the spare capacity of $c'$, which exists because $\sum c'_k\ge m$, at a further cost at most $W\|c-c'\|_1$. Exchange the roles of $c$ and $c'$. For the greedy description, let $y$ be feasible and suppose that $y_k>0$ while $y_j<c_j$ for some $j$ with $w_j<w_k$; moving the mass $\min(y_k,c_j-y_j)$ from $k$ to $j$ keeps $y$ feasible and does not increase the cost. Repeating this for the slots in order of increasing weight transforms any optimal solution, without increasing its cost, into the one that fills the slots in that order until the total $m$ is reached. For $F^-$, mass on a slot of nonnegative weight can be removed without increasing the cost, and the same argument applies to the slots of negative weight. If the capacities and $m$ are integers, the greedy solution is integral.
\end{proof}

\begin{lemma}[weak duality with one price]\label{lem:weak-duality}
With the notation of \cref{lem:lp-lipschitz}, for every $t\in\R$ and every $y$ with $0\le y_k\le c_k$ and $\sum y_k=m$,
\[
\sum_kw_ky_k\ \ge\ mt+\sum_kc_k\min(0,w_k-t).
\]
The same holds if only $\sum y_k\le m$ is assumed, provided $t\le0$.
\end{lemma}

\begin{proof}
$\sum_kw_ky_k=t\sum_ky_k+\sum_k(w_k-t)y_k$ and $(w_k-t)y_k\ge c_k\min(0,w_k-t)$. If $\sum y_k\le m$ and $t\le0$, then $t\sum_ky_k\ge tm$.
\end{proof}

\section{Numerical evidence for small \texorpdfstring{$n$}{n}}\label{app:numerics}

The polynomials $Q_n$ were computed exactly, from the definitions of \cref{sec:functional}, with programs written separately from those that produced the certificates, for $1\le n\le19$ ($\lambda=2/11$), $n\le11$ ($2/9$), $n\le19$ ($-2/9$) and $n\le16$ ($-2/7$). The entries $\Lambda_X(W_nx^{i+j})$ are affine in $X$; the determinant is obtained exactly by interpolation in $X$. Values at $\xi$ were enclosed rigorously in ball arithmetic. These computations do not enter the proof; they test the finite statements of \cref{sec:exceptional,sec:ordinary} and indicate how far the finite-$n$ data are from the asymptotic rates. They support the general arguments as follows.

\begin{enumerate}[(a)]
\item On sample entries for all four values, $\Lambda_\xi(W_nx^s)$ computed from the exact affine entries agrees with the series $\sum_{x\ge1}\lambda^x(xW_n(x)x^s)'$ to $40$ digits (\cref{lem:series}).
\item In every computed case $\deg Q_n=h$; moreover every computed $Q_n$ is irreducible over $\Q$, hence nonzero at every rational point. This is consistent with \cref{thm:nonvanishing} but is not used.
\item For every computed $n$ and every prime $p$, the actual Gauss valuation $\vg p(\Qt_n)$ satisfies the exact bounds of \cref{prop:generic,prop:two,prop:P,prop:Q} and, for ordinary $p$ with $\sqrt K<p\le K$, the bound $v_p(S_n^h/F_n)+\max(\cJ_p,\cP_p)$ of \cref{prop:jet,prop:pole}; and no prime $p>K$ outside $\Ex$ divides a denominator. The version of the pole bound without the zero-class correction of \cref{rem:zero-class} fails by exactly one unit at many pairs $(p,n)$, for instance for $\lambda=2/11$, $n=6$, at $p=19,23,29,31$.
\item The bounds at the exceptional primes are far from sharp in this range. For $\lambda=2/11$ the actual $\vg2(\Qt_n)$ exceeds the bound of \cref{prop:two} by $353$, $828$ and $2038$ for $n=5$, $10$ and $19$, that is, the bound underestimates the contribution of the prime $2$ to $\log\cont(\Qt_n)$ by about $0.61h^2$, $0.36h^2$ and $0.24h^2$. The relative slack decreases with $n$, and the data do not show whether the bound at $2$ is asymptotically sharp.
\end{enumerate}

\Cref{tab:numerics} lists the normalized content and value at the largest computed $n$, together with the thresholds for the asymptotic rates.

\begin{table}[ht]
\centering
\caption{Exact data at the largest computed $n$, normalized by $h^2$, and the thresholds $\rar$ (a lower bound for the lower limit of the content) and $\ran$ (an upper bound for the upper limit of the value). For $2/9$ the value at $n=11$ is above both. $P_n=\Qt_n/\cont(\Qt_n)$.}
\label{tab:numerics}
\begin{tabular}{@{}lcccccc@{}}
\toprule
$\lambda$ & $n$ & $h$ & $\frac{\log\cont\Qt_n}{h^2}$ & $\frac{\log|\Qt_n(\xi)|}{h^2}$ & $\frac{\log|P_n(\xi)|}{h^2}$ & $(\rar,\ \ran)$\\
\midrule
$2/11$ & $19$ & $76$ & $-0.1150$ & $-0.3481$ & $-0.2331$ & $(-0.2711,\ -0.3063)$\\
$2/9$ & $11$ & $88$ & $-1.1284$ & $-1.2886$ & $-0.1602$ & $(-1.33,\ -1.34)$\\
$-2/9$ & $19$ & $76$ & $-0.2393$ & $-0.6285$ & $-0.3891$ & $(-0.40,\ -0.49)$\\
$-2/7$ & $16$ & $80$ & $-0.3664$ & $-0.6563$ & $-0.2899$ & $(-0.5213,\ -0.5261)$\\
\bottomrule
\end{tabular}
\end{table}

These numbers are far from the asymptotic regime and are not evidence for the exact values of the rates. The normalized content is still decreasing in all four cases, as expected: the per-prime lower bounds converge to their limit only at the speed of the prime number theorem, and the contribution of the small primes $p\le\sqrt K$, bounded crudely in \cref{prop:generic}, decays like $n^{-1/2}$. The normalized value $h^{-2}\log|\Qt_n(\xi)|$ has stabilized for $2/11$, $-2/9$ and $-2/7$ well below the certified bounds, which suggests that the energy and contour estimates have room to spare in these cases. For $\lambda=2/9$ it is still decreasing, and at $n=11$ it lies above the bound $\ran$, and even above $\rar$, by about $0.05$: least-squares fits of the form $c_0+c_1\frac{\log h}h+\frac{c_2}h$ to the data for $n\ge3$, $5$ or $7$ give limits between $-1.349$ and $-1.341$, close to the certified bound $-1.3436$, while fits with more parameters are unstable. Since the comparison measure for $2/9$ is nearly optimal (\cref{sec:margins}), this is what one expects if the bound is close to the true limit, but it is equally what one would see if the true limit were slightly above the bound. The data therefore give no numerical corroboration for the case $2/9$, in contrast with the other three cases (\cref{sec:margins}).


% back matter starts on a new page
\clearpage
\afbacksection{Reproducibility and formal verification}

This section reports only verification that has actually been carried out. The status of each step is recorded in \cref{tab:status}: the general arguments are candidates, not independently verified, and the main theorems are not formalized. This preprint is contribution AF-0001, version 0.1.0, of the AxiomForge archive; paths below are relative to its directory \path{results/AF-0001-dilogarithm-irrationality} in \url{https://github.com/axiomforge-math/mathematics}. The record of every check, with its command, date, tool versions and outcome, is the file \path{logs/verification.md}. The programs are in \path{logs/experiments/scripts/}, their recorded outputs in \path{logs/experiments/outputs/}, and \path{logs/experiments/README.md} lists each program and what it establishes. They need Python~3.11 or later with \texttt{python-flint}~0.9.0 and \texttt{mpmath}~1.3.0 (\path{logs/experiments/requirements.txt}).

\emph{Certificate programs.} The finite computations of \cref{app:certificates} are carried out by the four original certificate programs \path{certify_two_elevenths.py}, \path{certify_two_ninths.py}, \path{certify_minus_two_ninths.py} and \path{certify_minus_two_sevenths.py}, which use only exact rational arithmetic and the Python standard library, together with \path{check_slot_duality.py} (the weak-duality minorants of \cref{sec:cert-profile} on all $5026$ subintervals) and \path{check_local_determinants.py} (the third computation of the values of \cref{tab:local-dets}). These programs signal failure through Python assertions and refuse to run when assertions are disabled. What they check is finite: the inequalities between the computed bounds and the thresholds of \cref{tab:rates}, and the identities of \cref{tab:local-dets}. The margins between the computed arithmetic and analytic bounds are $0.577$ in the $n^2$-scale for $2/11$, and $0.019104$, $0.11232$ and $0.0047556$ per $h^2$ for $2/9$, $-2/9$ and $-2/7$.

\emph{Separately written checkers.} The programs \path{verify_profile_integrals.py}, \path{verify_comparison_bounds.py} and \path{verify_rate_comparison.py} recompute the same quantities from the definitions in this paper rather than from the original programs, with exact rational arithmetic for the profile integrals and ball arithmetic for all transcendental quantities; each exits with an error status on a failed or undecided inequality. With the same comparison measures, exact profile integrals and cell enclosures on a mesh of width $1/6000$, they obtain the margins $1.084$ in the $n^2$-scale (that is, $0.0678$ per $h^2$) for $2/11$, and $0.019379$, $0.13978$ and $0.02969$ per $h^2$. The command \texttt{python3 run\_checks.py certificates}, run in \path{logs/experiments/scripts/}, executes the certificate programs and the checkers.

\emph{Exact computations for small $n$.} The program \path{compute_hankel_polynomials.py} computes the polynomials $Q_n$ of \cref{app:numerics} exactly; \path{check_series_identity.py}, \path{check_exceptional_primes.py}, \path{check_valuation_bounds.py}, \path{check_discrete_bounds.py}, \path{check_irreducibility.py} and \path{estimate_family_rates.py} produce the observations reported there. The stages \texttt{exact-data} and \texttt{exploratory} of \path{run_checks.py} run them. These computations are numerical evidence and do not enter the proof.

\emph{Lean.} The Lean~4 package in \path{lean/} (Lake package \texttt{axiomforge-af0001}, root module \texttt{AxiomForge.AF0001}, toolchain and Mathlib revision pinned in \path{lean/lean-toolchain} and \path{lean/lakefile.toml}) contains the definitions of the construction, the general irrationality criterion and the degree bound $\deg\Qt_n\le h$. It proves the eight local Gram determinants of \cref{tab:local-dets} by kernel-checked factorizations, whose Lean sources are generated by \path{generate_local_gram_certificates.py}, and the final logarithmic comparisons of the content rates for $2/11$ and $-2/7$. It states the main theorems only conditionally: the four theorems
\begin{center}\small
\texttt{irrational\_dilog\_two\_elevenths\_of\_estimates},\quad \texttt{irrational\_dilog\_two\_ninths\_of\_estimates},\\
\texttt{irrational\_dilog\_neg\_two\_ninths\_of\_estimates},\quad \texttt{irrational\_dilog\_neg\_two\_sevenths\_of\_estimates}
\end{center}
(namespace \texttt{AxiomForge.Dilogarithm}) derive the irrationality of $\Li_2(\lambda)$ from the hypotheses \texttt{HankelEstimates}, which state \cref{thm:arith,thm:arch,thm:nonvanishing} with the thresholds of \cref{tab:rates}. These hypotheses are not proved in Lean, and no unconditional formal statement of \cref{thm:main} exists; the formalization of the three estimates is in preparation. The build command is \texttt{lake exe cache get \&\& lake build}; the output of \texttt{\#print axioms} for the audited declarations is \path{lean/AXIOMS.txt}, and the outcome of the build and of the audit is recorded in \path{logs/verification.md}.

\emph{Review.} No independent human review of this work has taken place. The recomputations and checks above, and the adversarial review mentioned in \cref{sec:intro-status}, were carried out with AI systems: they are AI-internal checks, not an independent human or external verification.

\afbacksection{Acknowledgements}

We thank the authors of the works cited here for making their preprints and formal developments publicly available.

\afbacksection{AI-assisted research disclosure}

The original exploration, the constructions, the candidate proofs and the certificate programs underlying this paper were produced with GPT-6.2 XHigh (web version), as stated by the owners of AxiomForge, in the form of three separate candidate manuscripts with accompanying programs and an uncompiled Lean draft.

The reconstruction of the argument as one paper with a unified construction, the corrections it contains (in particular the zero-class correction of \cref{rem:zero-class}), the audit computations (\cref{app:numerics}), the recomputation of the certificates, this manuscript and the Lean formalization described in the preceding section were produced by Claude Opus 5.5 (Anthropic), working under the direction of AxiomForge.

The manuscript of 25 November 2026 was reviewed adversarially by Claude Fable 5.2 (Anthropic) in a separate session, and the present revision (attributions to earlier work, the complete proof of \cref{prop:profile-uniform}, corrections of statements and of the description of the verification record) was produced by the same model in response. That review and this revision are AI work; the review is not an independent review.

No independent human review of the mathematics has taken place. Cross-checks between AI systems, including the recomputations reported above, are not independent mathematical certification, and the involvement of AI systems in any part of this work is not evidence of its correctness.

\emph{Responsibility.} AI systems are not authors of this work and cannot take responsibility
for it. AxiomForge takes full responsibility for the entire content of this manuscript,
including every part produced with AI assistance, and for any errors it may contain. Reports
of errors are welcome at \href{mailto:axiomforge.math@gmail.com}{axiomforge.math@gmail.com};
confirmed errors are corrected openly in a new version, recorded in the changelog of this
contribution, with credit to whoever reported them.


\clearpage
\bibliographystyle{amsplain}
\bibliography{references}

\end{document}
